% Options for packages loaded elsewhere
\PassOptionsToPackage{unicode}{hyperref}
\PassOptionsToPackage{hyphens}{url}
\PassOptionsToPackage{dvipsnames,svgnames,x11names}{xcolor}
%
\documentclass[
  11pt,
  a4paper,
]{article}
\usepackage{amsmath,amssymb}
\usepackage{iftex}
\ifPDFTeX
  \usepackage[T1]{fontenc}
  \usepackage[utf8]{inputenc}
  \usepackage{textcomp} % provide euro and other symbols
\else % if luatex or xetex
  \usepackage{unicode-math} % this also loads fontspec
  \defaultfontfeatures{Scale=MatchLowercase}
  \defaultfontfeatures[\rmfamily]{Ligatures=TeX,Scale=1}
\fi
\usepackage{lmodern}
\ifPDFTeX\else
  % xetex/luatex font selection
\fi
% Use upquote if available, for straight quotes in verbatim environments
\IfFileExists{upquote.sty}{\usepackage{upquote}}{}
\IfFileExists{microtype.sty}{% use microtype if available
  \usepackage[]{microtype}
  \UseMicrotypeSet[protrusion]{basicmath} % disable protrusion for tt fonts
}{}
\makeatletter
\@ifundefined{KOMAClassName}{% if non-KOMA class
  \IfFileExists{parskip.sty}{%
    \usepackage{parskip}
  }{% else
    \setlength{\parindent}{0pt}
    \setlength{\parskip}{6pt plus 2pt minus 1pt}}
}{% if KOMA class
  \KOMAoptions{parskip=half}}
\makeatother
\usepackage{xcolor}
\usepackage{longtable,booktabs,array}
\usepackage{calc} % for calculating minipage widths
% Correct order of tables after \paragraph or \subparagraph
\usepackage{etoolbox}
\makeatletter
\patchcmd\longtable{\par}{\if@noskipsec\mbox{}\fi\par}{}{}
\makeatother
% Allow footnotes in longtable head/foot
\IfFileExists{footnotehyper.sty}{\usepackage{footnotehyper}}{\usepackage{footnote}}
\makesavenoteenv{longtable}
\setlength{\emergencystretch}{3em} % prevent overfull lines
\providecommand{\tightlist}{%
  \setlength{\itemsep}{0pt}\setlength{\parskip}{0pt}}
\setcounter{secnumdepth}{-\maxdimen} % remove section numbering
% Layout of the GZero documents (used by docs/build_manual.py through pandoc).
\usepackage[a4paper,top=25mm,bottom=25mm,left=25mm,right=25mm,headheight=14pt]{geometry}
\usepackage{lmodern}
\renewcommand{\familydefault}{\sfdefault}
\usepackage[T1]{fontenc}
\usepackage{microtype}
\usepackage{amsmath,amssymb}
\usepackage{xcolor}
\definecolor{gzblue}{HTML}{0072B2}
\definecolor{gzgrey}{HTML}{F3F4F6}
\usepackage{titlesec}
\titleformat{\section}{\Large\bfseries\color{gzblue}}{}{0pt}{}[\vspace{2pt}\titlerule]
\titleformat{\subsection}{\large\bfseries}{}{0pt}{}
\titleformat{\subsubsection}{\normalsize\bfseries}{}{0pt}{}
\titlespacing*{\section}{0pt}{18pt plus 4pt minus 2pt}{8pt}
\titlespacing*{\subsection}{0pt}{12pt plus 3pt minus 2pt}{4pt}
\usepackage{fancyhdr}
\pagestyle{fancy}
\fancyhf{}
\makeatletter\AtBeginDocument{\fancyhead[L]{\small\@title}}\makeatother
\fancyhead[R]{\small\nouppercase{\rightmark}}
\fancyfoot[C]{\small\thepage}
\renewcommand{\headrulewidth}{0.4pt}
\renewcommand{\sectionmark}[1]{\markboth{#1}{#1}}
\renewcommand{\subsectionmark}[1]{}
\fancypagestyle{plain}{\fancyhf{}\fancyfoot[C]{\small\thepage}\renewcommand{\headrulewidth}{0pt}}
\setlength{\parindent}{0pt}
\setlength{\parskip}{5pt plus 1pt}
\usepackage{etoolbox}
% tables: smaller type, a little more row height, ragged cells
\AtBeginEnvironment{longtable}{\small\renewcommand{\arraystretch}{1.25}}
\usepackage{ragged2e}
\setlength{\LTpre}{6pt}
\setlength{\LTpost}{6pt}
\usepackage{enumitem}
\setlist{itemsep=2pt,topsep=4pt}
% code blocks on a light background
\usepackage{fvextra}
\DefineVerbatimEnvironment{Highlighting}{Verbatim}{breaklines,commandchars=\\\{\}}
\BeforeBeginEnvironment{verbatim}{\begingroup\small}
\AfterEndEnvironment{verbatim}{\endgroup}
\usepackage[most]{tcolorbox}
\usepackage{tocloft}
\renewcommand{\cftsecfont}{\bfseries}
\renewcommand{\cftsecpagefont}{\bfseries}
\setlength{\cftbeforesecskip}{3pt}
\setcounter{tocdepth}{2}
\usepackage{hyperref}
\hypersetup{colorlinks=true,linkcolor=gzblue,urlcolor=gzblue,citecolor=gzblue,pdfauthor={arman-rashid}}
\urlstyle{same}
\sloppy
\ifLuaTeX
  \usepackage{selnolig}  % disable illegal ligatures
\fi
\IfFileExists{bookmark.sty}{\usepackage{bookmark}}{\usepackage{hyperref}}
\IfFileExists{xurl.sty}{\usepackage{xurl}}{} % add URL line breaks if available
\urlstyle{same}
\hypersetup{
  pdftitle={GZero - User Manual},
  colorlinks=true,
  linkcolor={Maroon},
  filecolor={Maroon},
  citecolor={Blue},
  urlcolor={Blue},
  pdfcreator={LaTeX via pandoc}}

\title{GZero - User Manual}
\usepackage{etoolbox}
\makeatletter
\providecommand{\subtitle}[1]{% add subtitle to \maketitle
  \apptocmd{\@title}{\par {\large #1 \par}}{}{}
}
\makeatother
\subtitle{Break-junction analysis software. Version 1.4, October 2026}
\author{}
\date{}

\begin{document}
\maketitle

\renewcommand*\contentsname{Contents}
{
\hypersetup{linkcolor=}
\setcounter{tocdepth}{3}
\tableofcontents
}
\hypertarget{introduction-and-installation}{%
\section{1. Introduction and
installation}\label{introduction-and-installation}}

I wrote GZero to analyse single-molecule break-junction data (MCBJ and
STM-BJ) in one place, from the raw TDMS files of the setup to the
figures for a paper. It converts the raw signals to conductance, cuts
out the single traces, makes the usual histograms, measures plateau
lengths, calculates 2D correlation maps, clusters traces and gets the
flicker-noise exponent. It also tells you how often a molecule actually
bridged the gap (the junction yield) and whether the conductance drifted
during the run. For I-V sweeps it fits the single-level model and finds
the transition voltage. For conductance-time traces it finds flickering
and mechanical switching events, a separate tab does the ensemble
statistics of large-area (EGaIn) junctions, and another one the
thermopower (Seebeck coefficient) from thermovoltage measurements. Piezo
modulation is covered too, and there is a simulator that makes fake data
with known answers, which I use to check settings.

\hypertarget{the-tabs}{%
\subsection{1.1 The tabs}\label{the-tabs}}

\begin{longtable}[]{@{}
  >{\raggedright\arraybackslash}p{(\columnwidth - 2\tabcolsep) * \real{0.2795}}
  >{\raggedright\arraybackslash}p{(\columnwidth - 2\tabcolsep) * \real{0.6905}}@{}}
\toprule\noalign{}
\begin{minipage}[b]{\linewidth}\raggedright
Tab
\end{minipage} & \begin{minipage}[b]{\linewidth}\raggedright
What it's for
\end{minipage} \\
\midrule\noalign{}
\endhead
\bottomrule\noalign{}
\endlastfoot
Convert TDMS & Raw TDMS channels to conductance, \texttt{.cvr} export \\
Data \& Traces & Loading files, finding and cutting traces, looking
through them, marking bad ones \\
Histograms & 1D log G histogram, 2D conductance-distance histogram,
Gaussian fits \\
Plateau length & How long the molecular plateau is, in up to three
windows; tunnelling decay \\
Junction statistics & Which traces have a molecular plateau (junction
yield), the conductance of each plateau, drifts over the measurement \\
Correlation & 2D cross-correlation of conductance values within
traces \\
Clustering & PCA and clustering of traces into groups \\
Flicker noise & Noise power against conductance and the scaling exponent
n \\
Events & Flickering and mechanical switching in G(t), conductance-time
histograms of the holds \\
I-V & Forward and backward bias sweeps, dI/dV, 2D I-V histograms,
single-level model fit, transition voltage spectroscopy \\
Piezo modulation & Decay constant $\beta$ from a modulated piezo \\
EGaIn (ensemble) & Large-area junctions: log\textbar J\textbar{}
statistics, yield, rectification, $V_\mathrm{trans}$, $\beta$ from a
length series \\
Thermopower & Thermovoltage histograms per temperature difference,
Seebeck coefficient, S against conductance \\
Simulation & Synthetic recordings \\
\end{longtable}

\hypertarget{what-you-need}{%
\subsection{1.2 What you need}\label{what-you-need}}

Windows 10 or 11 (64 bit) for the exe. From Python it should also run on
macOS and Linux, though I've only used it on Windows. 8 GB of RAM is
fine for normal files. Big files are loaded completely into memory: my 1
GB test file (73 min at 10 kHz) took about 7 s to load, and I'd plan
roughly 3-4 GB of RAM per GB of TDMS. The window needs about 1400 x 900
pixels to be comfortable.

\hypertarget{installing-and-starting}{%
\subsection{1.3 Installing and starting}\label{installing-and-starting}}

The simplest way is the exe. Copy the whole \texttt{GZero} folder
somewhere,
e.g.~\texttt{C:\textbackslash{}Programs\textbackslash{}GZero}, and
double-click \texttt{GZero.exe}. Don't move the exe out of the folder,
it needs the files next to it. This manual is in
\texttt{docs\textbackslash{}User\_\allowbreak{}Manual.pdf} and also
opens from Help -\textgreater{} User manual (F1).

The source code is on GitHub: https://github.com/arman-rashid/GZero
(Apache License 2.0). To run it from Python (3.9 or newer), go to the
project folder (the one with \texttt{GZero.py}) and run

\begin{verbatim}
py -m pip install -r requirements.txt
py GZero.py
\end{verbatim}

It needs PySide6, numpy, scipy, matplotlib, pandas, scikit-learn, npTDMS
and igor2 (for Igor \texttt{.pxp} / \texttt{.ibw} files).

If you change the code and want a new exe:
\texttt{py -m pip install pyinstaller}, then
\texttt{py build\_\allowbreak{}exe.py} in the same folder. The new
version ends up in \texttt{dist\textbackslash{}GZero\textbackslash{}}.

\hypertarget{the-window}{%
\subsection{1.4 The window}\label{the-window}}

Every tab has the settings and buttons on the left and the plots on the
right. The toolbar above each plot lets you zoom, pan and save what's on
screen. When you save results, each plot is also written as a separate
figure in journal format (13.1). Most settings have a tooltip, so hover
over them if something isn't clear. The status bar at the bottom says
what the program is doing; longer calculations run in the background and
the window doesn't freeze.

Settings are only kept while the program is open. File -\textgreater{}
Save settings writes them to a file and File -\textgreater{} Load
settings reads them back. You can also load the
\texttt{*\_\allowbreak{}config.json} from an old results folder, which
is the easiest way to repeat an analysis with exactly the same settings.

\hypertarget{quick-start}{%
\section{2. Quick start}\label{quick-start}}

This is the short version of a normal analysis, from a raw file to a
histogram and a noise exponent. The chapter numbers point to the
details.

\begin{enumerate}
\def\labelenumi{\arabic{enumi}.}
\tightlist
\item
  Go to Data \& Traces, click Load files and pick a \texttt{.tdms},
  \texttt{.cvr}, text or Igor \texttt{.pxp} file (several at once is
  fine, or use Load folder). For everything except TDMS check the column
  settings first: which column is what, and the bias if the file has no
  voltage column (4.2).
\item
  The traces are found automatically after the first load. Under the
  button you see how many were found, how many were accepted and why the
  rest were thrown out.
\item
  Click through a few traces with \textless{} and \textgreater{} or the
  arrow keys. Green means accepted, red rejected. If good-looking traces
  are being rejected, change the limits (4.4) and press Detect / cut
  traces again.
\item
  G and B mark the current trace as good or bad by hand.
\item
  Open Histograms. The 1D and 2D histograms are already there.
\item
  Put \emph{fit from} and \emph{fit to} around the molecular peak,
  choose how many Gaussians, click Fit peaks (5).
\item
  In Plateau length, set window 1 around the peak and click Calculate
  (6).
\item
  In Flicker noise, click the Morris et al.~2025 preset, then Peak from
  histogram fit, then Run noise analysis. n and its error appear under
  the button (9).
\item
  Correlation and Clustering if you need them (7, 8). Junction
  statistics gives the yield and checks for drift (6.4).
\item
  File -\textgreater{} Save results (Ctrl+S), choose where and give the
  folder a name. Everything goes in there: text files, figures, settings
  (13).
\end{enumerate}

If you just want to try the program, go to Simulation and click
Generate. You get a fake data set where you know the right answers, and
you can see if the analysis finds them (12).

\hypertarget{data-formats-and-the-convert-tab}{%
\section{3. Data formats and the Convert
tab}\label{data-formats-and-the-convert-tab}}

You only need the Convert tab for raw \texttt{.tdms} files. It turns
them into conductance (G/$G_0$) and either writes \texttt{.cvr} files or
loads the result straight into the analysis. All other file types are
loaded in Data \& Traces (4).

\hypertarget{file-types}{%
\subsection{3.1 File types}\label{file-types}}

\begin{longtable}[]{@{}
  >{\raggedright\arraybackslash}p{(\columnwidth - 4\tabcolsep) * \real{0.2801}}
  >{\raggedright\arraybackslash}p{(\columnwidth - 4\tabcolsep) * \real{0.4882}}
  >{\raggedright\arraybackslash}p{(\columnwidth - 4\tabcolsep) * \real{0.2017}}@{}}
\toprule\noalign{}
\begin{minipage}[b]{\linewidth}\raggedright
Format
\end{minipage} & \begin{minipage}[b]{\linewidth}\raggedright
What's in it
\end{minipage} & \begin{minipage}[b]{\linewidth}\raggedright
Where to load it
\end{minipage} \\
\midrule\noalign{}
\endhead
\bottomrule\noalign{}
\endlastfoot
\texttt{.tdms}, raw & Group \texttt{BJRawData} with the bias
(\texttt{M\_\allowbreak{}E\#} or \texttt{M\_\allowbreak{}E1-E2}),
low-gain current \texttt{M\_\allowbreak{}i\#a}, high-gain current
\texttt{M\_\allowbreak{}i\#b} and piezo \texttt{U\_\allowbreak{}Piezo},
plus the \texttt{Setup}, \texttt{Calibration\_\allowbreak{}M\#} and
\texttt{Events} groups & Convert tab or Data \& Traces \\
\texttt{.tdms}, with G & A \texttt{G} channel in siemens and
\texttt{U\_\allowbreak{}Piezo} (the \texttt{*\_\allowbreak{}M.tdms}
files) & Data \& Traces \\
\texttt{.tdms}, other & Any channel names: type them into the channel
settings (3.3) & Convert tab or Data \& Traces \\
\texttt{.cvr} & Two tab-separated columns, back-calculated current (A)
and piezo (nm). Old files with three columns work too & Data \&
Traces \\
\texttt{.txt} \texttt{.csv} \texttt{.dat} \texttt{.tsv} & Any number
columns; you say which column is what & Data \& Traces \\
\texttt{.pxp} (Igor packed experiment), \texttt{.ibw} (Igor binary wave)
& Every numeric wave is a column, named after the wave
(\texttt{folder:wave} inside data folders, \texttt{wave[0]},
\texttt{wave[1]} \ldots{} for 2D waves) & Data \& Traces \\
\end{longtable}

TDMS files carry the sampling rate, channel gains, piezo scale and
series resistor. Text files don't, so for those you type them into the
column settings (4.2). Igor waves carry their x scaling: if a wave has x
in seconds, the sampling rate comes from it.

The program never recalculates what a file already gives. If there is a
conductance column or channel, G is taken from it. If there is a voltage
column or bias channel, it is used as the bias and the constant bias is
ignored. Piezo, current and time columns are also used as they are. Only
what's missing is calculated: G from the current and the bias, or the
current from G and the bias.

The current can come from a single-stage amplifier (one current, Ia) or
a dual-stage amplifier (a low-gain stage Ia and a high-gain stage Ib
recorded at the same time). The \emph{current amplifier} setting says
which one (3.3 for TDMS, 4.2 for the other files). With a single stage
there is nothing to merge and all the Ib settings are greyed out.

\hypertarget{using-the-tab}{%
\subsection{3.2 Using the tab}\label{using-the-tab}}

\begin{enumerate}
\def\labelenumi{\arabic{enumi}.}
\tightlist
\item
  Add files or Add folder. When you click a file, the panel shows its
  sampling rate, length, modules, channel gains, series resistor and
  whether diode calibration data is in it.
\item
  Change the conversion settings if needed (3.3). Empty fields mean
  ``take it from the file''.
\item
  Preview selected shows the raw current channels, log(G/$G_0$) with the
  event marks (green = approach, red = retract, grey = hold) and the
  piezo.
\item
  Then either Convert all -\textgreater{} .cvr, which writes one
  \texttt{.cvr} per TDMS file (tick Also write .csv files if you also
  want a table with time, G, log G, piezo, bias and current), or Load
  all into analysis, which converts in memory and jumps to Data \&
  Traces.
\end{enumerate}

\hypertarget{conversion-settings}{%
\subsection{3.3 Conversion settings}\label{conversion-settings}}

\begin{longtable}[]{@{}
  >{\raggedright\arraybackslash}p{(\columnwidth - 4\tabcolsep) * \real{0.2186}}
  >{\raggedright\arraybackslash}p{(\columnwidth - 4\tabcolsep) * \real{0.5456}}
  >{\raggedright\arraybackslash}p{(\columnwidth - 4\tabcolsep) * \real{0.2058}}@{}}
\toprule\noalign{}
\begin{minipage}[b]{\linewidth}\raggedright
Setting
\end{minipage} & \begin{minipage}[b]{\linewidth}\raggedright
Meaning
\end{minipage} & \begin{minipage}[b]{\linewidth}\raggedright
Default
\end{minipage} \\
\midrule\noalign{}
\endhead
\bottomrule\noalign{}
\endlastfoot
module & Which measurement module to use (\texttt{M\_\allowbreak{}E\#},
\texttt{M\_\allowbreak{}i\#a}, \texttt{M\_\allowbreak{}i\#b}) & 1 \\
current amplifier & \texttt{dual stage}: Ia (low gain) and Ib (high
gain) are merged. \texttt{single stage}: only the Ia channel is used,
nothing is merged & dual stage \\
Ia channel / Ib channel & \texttt{auto} =
\texttt{M\_\allowbreak{}i\textless{}module\textgreater{}a} /
\texttt{M\_\allowbreak{}i\textless{}module\textgreater{}b}. Type another
channel name if your file uses other names; a channel that is already in
A needs gain 1 & auto \\
conductance channel & \texttt{auto} uses a \texttt{G} channel only if
there are no current channels. Type a channel name to take G directly
from that channel, \texttt{none} to always calculate it & auto \\
conductance unit & Unit of the conductance channel: \texttt{auto} (from
the channel, otherwise S), \texttt{S}, \texttt{G/\allowbreak{}G0} or
\texttt{log G/\allowbreak{}G0} & auto \\
piezo channel & \texttt{auto} = \texttt{U\_\allowbreak{}Piezo}; another
name, or \texttt{none} & auto \\
calibration & \texttt{linear}: I = (V - offset) x gain x correction.
\texttt{diode}: the piecewise polynomial stored in
\texttt{Calibration\_\allowbreak{}M\#} (log amplifier) & linear \\
gain Ia / gain Ib (A/V) & Amplifier gains; empty = the
\texttt{scale\_\allowbreak{}factor} stored with the channel & from
file \\
offset / correction Ia, Ib & Offset (V) and correction factor for the
linear calibration & 0 / 1 \\
amplification & Extra factor on the current & 1 \\
invert current & Flips the sign & off \\
combine Ia / Ib & Dual stage only. \texttt{auto} uses the high-gain Ib
unless it's saturated, then Ia. \texttt{a} or \texttt{b} uses only one &
auto \\
Ib saturation (V) & Above this, Ib counts as saturated & 10 \\
Ia saturation (V) & Above this, the contact conducts more than the
amplifier can measure and the sample gets the saturated G & 10.4 \\
switch to Ib below log G & Optional: only use Ib below this log(G/$G_0$)
& off \\
bias / bias channel & \texttt{measured} takes the bias channel as an
array (\texttt{auto} = \texttt{M\_\allowbreak{}E1-E2} if it exists,
otherwise \texttt{M\_\allowbreak{}E\textless{}module\textgreater{}}, or
type a name) and the fixed bias isn't used; \texttt{fixed} uses the
fixed bias. Without a bias channel the fixed bias is used & measured,
auto \\
series R compensation / series R & Corrects for the current-limiting
resistor, G = I / (E - I x R). Empty R = value from the Setup group &
on, from file \\
saturated G ($G_0$) & What saturated samples are set to & 100 \\
piezo scale (nm/V) & Piezo voltage to elongation; empty = the channel's
\texttt{scale\_\allowbreak{}factor} & from file \\
bias / current / conductance filter & none, Savitzky-Golay, median or
moving average. The window is 2 x side points + 1. Can work on log
values & bias: moving average, 25 points \\
\end{longtable}

\hypertarget{what-happens-in-the-conversion}{%
\subsection{3.4 What happens in the
conversion}\label{what-happens-in-the-conversion}}

The current is calibrated and, for a dual-stage amplifier, the two gain
channels are put together (Ib where it isn't saturated, Ia elsewhere).
The conductance is the current divided by the voltage that's actually
across the junction, i.e.~the bias minus what drops over the series
resistor:

\[
G = \frac{I}{E - I R} \cdot \frac{1}{G_0}, \qquad G_0 = \frac{2e^2}{h} = 7.748 \times 10^{-5}\ \mathrm{S}
\]

If the file has a conductance channel and you choose it, none of this is
done: G comes straight from the channel and only the piezo and bias are
read. If the resistor takes the whole bias (E - I x R \textless= 0), or
the low-gain amplifier saturates, the junction is in contact and beyond
what can be measured, so the sample is set to the saturated G.

The \texttt{.cvr} file contains the back-calculated current, I\_bc = G x
$G_0$ x E, which is the current you'd get without the series resistor.
When a \texttt{.cvr} is loaded, G is calculated back as I\_bc / bias /
$G_0$, so the bias in the column settings has to be right.

About the diode calibration: \texttt{Calibration\_\allowbreak{}M\#}
holds five polynomials (PH, PL, M, NL, NH, coefficients in ascending
order), an offset and two thresholds. With u = V - offset, the middle
polynomial gives I directly when \textbar u\textbar{} is below the low
threshold; outside that, the low or high polynomial gives log10
\textbar I\textbar, and negative u uses the negative-side polynomials
with a minus sign. I haven't checked this against a measurement with
module 2 yet, so test it with a known resistor before trusting it.

\hypertarget{data-traces-tab}{%
\section{4. Data \& Traces tab}\label{data-traces-tab}}

Here the recordings are loaded and cut into single breaking traces
(making traces too, if you want them). All the other tabs only see the
traces that are accepted here.

\hypertarget{loading}{%
\subsection{4.1 Loading}\label{loading}}

Load files takes one or several files, Load folder takes every
\texttt{.tdms}, \texttt{.cvr}, \texttt{.txt}, \texttt{.csv},
\texttt{.dat}, \texttt{.tsv}, \texttt{.pxp} and \texttt{.ibw} in a
folder. Everything you load is analysed together. TDMS files are
converted with the settings from the Convert tab (3.3), text files with
the column settings below. Remove drops the selected recordings and
their traces, Clear drops everything. Show columns lists the columns (or
Igor waves) of a file with their numbers and lengths, so you know what
to type into the column settings. The tooltip of a loaded recording says
which quantities came from the file and which were calculated. The list
shows how long each recording is and whether it has a piezo signal;
click one to see it in the overview plot.

\hypertarget{column-settings-text-.cvr-and-igor-files}{%
\subsection{4.2 Column settings (text, .cvr and Igor
files)}\label{column-settings-text-.cvr-and-igor-files}}

A column is given by its number (0 is the first) or by its name: the
header of a text file or the wave name in an Igor file. Empty means the
file doesn't have it.

\begin{longtable}[]{@{}
  >{\raggedright\arraybackslash}p{(\columnwidth - 4\tabcolsep) * \real{0.2752}}
  >{\raggedright\arraybackslash}p{(\columnwidth - 4\tabcolsep) * \real{0.5262}}
  >{\raggedright\arraybackslash}p{(\columnwidth - 4\tabcolsep) * \real{0.1685}}@{}}
\toprule\noalign{}
\begin{minipage}[b]{\linewidth}\raggedright
Setting
\end{minipage} & \begin{minipage}[b]{\linewidth}\raggedright
Meaning
\end{minipage} & \begin{minipage}[b]{\linewidth}\raggedright
Default
\end{minipage} \\
\midrule\noalign{}
\endhead
\bottomrule\noalign{}
\endlastfoot
conductance column & G given directly. If set, G is read from it and not
calculated & empty \\
conductance unit & \texttt{G/\allowbreak{}G0}, \texttt{S} or
\texttt{log G/\allowbreak{}G0} & G/$G_0$ \\
current amplifier & \texttt{single stage}: one current column.
\texttt{dual stage}: Ia and Ib, merged & single stage \\
current Ia column & The current (single stage) or the low-gain stage Ia
(dual stage) & 0 \\
current Ib column & The high-gain stage Ib (dual stage only) & empty \\
Ia scale / Ib scale (A per unit) & Turns the column into A: 1 if it is
in A, the amplifier gain in A/V if it is in V & 1 \\
Ib saturated from \textbar Ib\textbar{} & Column value from which Ib
counts as saturated and Ia is used instead; empty = 99.5 \% of the
largest \textbar Ib\textbar{} in the file & empty \\
Ia saturated from \textbar Ia\textbar{} & Column value from which the
contact is beyond the range, G = saturated G & off \\
use Ib below log G & Optional: only use Ib below this log(G/$G_0$) &
off \\
voltage column & Measured bias in V. Replaces the constant bias, and I-V
analysis needs it & empty \\
constant bias (V) & Only used without a voltage column. For
\texttt{.cvr} files this must be the bias of the measurement & 0.1 \\
series R (ohm) & Series resistor, G = I / (V - I x R) / $G_0$ & 0 \\
saturated G ($G_0$) & What saturated samples are set to & 100 \\
piezo / distance column & Piezo position or distance; empty if there is
none (then distance = point number x piezo rate, 4.5) & 1 \\
piezo scale (to nm) & Factor that turns the piezo column into nm & 1 \\
time column (s) & If there is one, the sampling rate is taken from it &
empty \\
sampling rate (Hz) & Used when neither a time column nor the Igor x
scaling gives it & 10000 \\
\end{longtable}

What is calculated depends on what you give:

\begin{longtable}[]{@{}
  >{\raggedright\arraybackslash}p{(\columnwidth - 4\tabcolsep) * \real{0.2984}}
  >{\raggedright\arraybackslash}p{(\columnwidth - 4\tabcolsep) * \real{0.2984}}
  >{\raggedright\arraybackslash}p{(\columnwidth - 4\tabcolsep) * \real{0.3733}}@{}}
\toprule\noalign{}
\begin{minipage}[b]{\linewidth}\raggedright
Columns given
\end{minipage} & \begin{minipage}[b]{\linewidth}\raggedright
Conductance
\end{minipage} & \begin{minipage}[b]{\linewidth}\raggedright
Current
\end{minipage} \\
\midrule\noalign{}
\endhead
\bottomrule\noalign{}
\endlastfoot
conductance (+ current) & from the file & from the file, or G x $G_0$ x
V \\
Ia (single stage) & I / (V - I R) / $G_0$ & Ia \\
Ia and Ib (dual stage) & from the merged current & Ib where not
saturated, Ia elsewhere \\
\end{longtable}

V is the voltage column if there is one, otherwise the constant bias.

Files with one wave or column per trace (common in Igor): put a
\texttt{*} in the conductance or Ia column,
e.g.~\texttt{G\_\allowbreak{}*}. Every matching column becomes its own
recording. Other columns with a \texttt{*} are paired by the same text,
so with piezo column \texttt{Z\_\allowbreak{}*} the wave
\texttt{G\_\allowbreak{}12} gets \texttt{Z\_\allowbreak{}12}. Columns of
different length (Igor waves) are cut to the shortest one.

Header lines are skipped, and tab, comma, semicolon or space as
separator all work.

\hypertarget{how-the-traces-are-found}{%
\subsection{4.3 How the traces are
found}\label{how-the-traces-are-found}}

With a piezo signal (\texttt{segmentation = piezo}) the piezo is
smoothed and its speed calculated. Each point is then either moving
forward, moving back, or held, which means the speed is below \emph{hold
if speed below (fraction)} times the typical speed. Very short pieces
(below \emph{min segment}) are merged with their neighbours. Every
moving piece is a candidate trace and every held piece is a hold.

A moving piece counts as breaking if the conductance drops by more than
a decade from start to end, and as making if it rises. If you'd rather
decide by piezo direction, set \emph{breaking direction} to
\texttt{piezo+} or \texttt{piezo-}.

A lot of measurements stop the piezo in the middle of a pull to record
G(t) for the noise analysis. With \emph{join pull-hold-pull} switched on
(default) the pulls before and after the stop are joined into one trace,
but the hold points themselves are left out, otherwise the time spent
standing still would show up as a fake peak in the histogram. The hold
is kept separately for the noise tab (9).

Without a piezo signal (\texttt{segmentation = threshold}) the contacts
are found where log G goes above the high limit, and each gap between
two contacts is split at its lowest point into a breaking and a making
trace.

\hypertarget{cutting-limits-and-checks}{%
\subsection{4.4 Cutting limits and
checks}\label{cutting-limits-and-checks}}

All limits are log10(G/$G_0$). Making traces are turned around and cut
the same way, so every trace goes from contact to tunnelling.

\begin{longtable}[]{@{}
  >{\raggedright\arraybackslash}p{(\columnwidth - 4\tabcolsep) * \real{0.2465}}
  >{\raggedright\arraybackslash}p{(\columnwidth - 4\tabcolsep) * \real{0.5275}}
  >{\raggedright\arraybackslash}p{(\columnwidth - 4\tabcolsep) * \real{0.1959}}@{}}
\toprule\noalign{}
\begin{minipage}[b]{\linewidth}\raggedright
Setting
\end{minipage} & \begin{minipage}[b]{\linewidth}\raggedright
Meaning
\end{minipage} & \begin{minipage}[b]{\linewidth}\raggedright
Default
\end{minipage} \\
\midrule\noalign{}
\endhead
\bottomrule\noalign{}
\endlastfoot
high limit to check & The trace has to get above this (metal contact),
otherwise it's rejected & 0 \\
low limit to check & The trace has to get below this, otherwise it's
rejected & -4.5 \\
start at & The cut trace starts at the last point above this; empty =
the high limit & empty \\
low limit to cut & The trace ends at the first point below this after
the contact \ldots{} & -5.0 \\
additional length & \ldots{} plus this many points & 100 \\
reject noisy start & Throws out traces whose first \emph{noise points}
scatter more than \emph{noise limit} (std of log G) & off, 25, 0.2 \\
max floor noise & Optional, throws out traces whose last points are
noisier than this & off \\
use traces & \texttt{break}, \texttt{make} or \texttt{both} & break \\
\end{longtable}

I usually put the low limits about half a decade above the noise floor.
You can see the floor in the overview plot as the flat band at the
bottom.

\hypertarget{distance-axis}{%
\subsection{4.5 Distance axis}\label{distance-axis}}

\begin{longtable}[]{@{}
  >{\raggedright\arraybackslash}p{(\columnwidth - 4\tabcolsep) * \real{0.2618}}
  >{\raggedright\arraybackslash}p{(\columnwidth - 4\tabcolsep) * \real{0.5450}}
  >{\raggedright\arraybackslash}p{(\columnwidth - 4\tabcolsep) * \real{0.1633}}@{}}
\toprule\noalign{}
\begin{minipage}[b]{\linewidth}\raggedright
Setting
\end{minipage} & \begin{minipage}[b]{\linewidth}\raggedright
Meaning
\end{minipage} & \begin{minipage}[b]{\linewidth}\raggedright
Default
\end{minipage} \\
\midrule\noalign{}
\endhead
\bottomrule\noalign{}
\endlastfoot
distance from & \texttt{piezo} uses the smoothed piezo signal,
\texttt{index} uses point number x piezo rate & piezo \\
displacement ratio & Junction displacement per nm of piezo (the MCBJ
attenuation factor; 1 for STM) & 1 \\
piezo rate & nm per point, for \texttt{index} and for files without
piezo & 0.001 \\
zero set & Distance zero is where the trace first drops below this log G
(-0.3 is 0.5 $G_0$) & -0.3 \\
piezo smoothing & Smoothing window for direction detection and the
distance axis & 10 ms \\
\end{longtable}

The piezo has to be smoothed because in my files its noise (about 0.003
nm per sample) is as big as the step between samples. For an MCBJ the
displacement ratio can be calibrated with the tunnelling fit in the
Plateau length tab (6.3).

\hypertarget{overview-and-trace-browser}{%
\subsection{4.6 Overview and trace
browser}\label{overview-and-trace-browser}}

The overview plot shows log G against time for the selected recording,
with the piezo on the right axis. Breaking traces are shaded red, making
traces blue and holds grey; light shading means rejected.

Below it, the trace browser shows one trace against distance with its
own histogram next to it. The line above the plot gives the trace
number, type, file, length and whether it's accepted (and if not, why).

\begin{longtable}[]{@{}
  >{\raggedright\arraybackslash}p{(\columnwidth - 2\tabcolsep) * \real{0.3920}}
  >{\raggedright\arraybackslash}p{(\columnwidth - 2\tabcolsep) * \real{0.5780}}@{}}
\toprule\noalign{}
\begin{minipage}[b]{\linewidth}\raggedright
Key / button
\end{minipage} & \begin{minipage}[b]{\linewidth}\raggedright
What it does
\end{minipage} \\
\midrule\noalign{}
\endhead
\bottomrule\noalign{}
\endlastfoot
Left / Right arrow, \textless{} \textgreater{} & Previous / next
trace \\
G, Good & Mark as good and go to the next one \\
B, Bad & Mark as bad and go to the next one \\
Auto & Go back to the automatic decision \\
Accepted only & Only browse accepted traces \\
Overlay & Plot N random accepted traces on top of each other \\
\end{longtable}

Manual marks are kept until you detect the traces again.

\hypertarget{histograms-tab}{%
\section{5. Histograms tab}\label{histograms-tab}}

The 1D histogram of log(G/$G_0$) over all accepted traces is where the
molecular conductance shows up as a peak. The 2D histogram (log G
against distance) shows how far junctions get stretched at each
conductance. Both are redrawn automatically whenever the traces change.

\hypertarget{settings}{%
\subsection{5.1 Settings}\label{settings}}

\begin{longtable}[]{@{}
  >{\raggedright\arraybackslash}p{(\columnwidth - 4\tabcolsep) * \real{0.2572}}
  >{\raggedright\arraybackslash}p{(\columnwidth - 4\tabcolsep) * \real{0.5336}}
  >{\raggedright\arraybackslash}p{(\columnwidth - 4\tabcolsep) * \real{0.1791}}@{}}
\toprule\noalign{}
\begin{minipage}[b]{\linewidth}\raggedright
Setting
\end{minipage} & \begin{minipage}[b]{\linewidth}\raggedright
Meaning
\end{minipage} & \begin{minipage}[b]{\linewidth}\raggedright
Default
\end{minipage} \\
\midrule\noalign{}
\endhead
\bottomrule\noalign{}
\endlastfoot
traces & All accepted traces, or one cluster from the Clustering tab &
all \\
log G min / max & Conductance range of both histograms & -6.5 / 1 \\
bins / decade & Conductance resolution & 100 \\
distance min / max, distance bins & Range and resolution of the 2D
histogram & -0.5 / 2 nm, 200 \\
normalization & \texttt{per trace} divides by the number of traces, so
different data sets can be compared. \texttt{counts} is raw counts,
\texttt{density} has area 1 & per trace \\
2D colormap, colour max & Colour scale of the 2D histogram. The top of
the scale is a percentile of the filled bins, so a few very full bins
don't wash everything out & viridis, 99.5 \\
overlay clusters in 1D & If there are clusters and \emph{all} is
selected, each cluster's histogram is drawn too, scaled by its share of
the traces & on \\
\end{longtable}

\hypertarget{fitting-the-peak}{%
\subsection{5.2 Fitting the peak}\label{fitting-the-peak}}

Set \emph{fit from} and \emph{fit to} around the peak (keep away from
the $G_0$ peak and the noise floor), choose the number of Gaussians and
whether you want a constant baseline, and click Fit peaks. The fit is a
sum of Gaussians in log G plus the optional constant:

\[
N(x) = c + \sum_k A_k \exp\left(-\frac{(x - \mu_k)^2}{2\sigma_k^2}\right), \qquad x = \log_{10}(G/G_0)
\]

The starting values are the highest maxima in the range. For every peak
you get the centre with its error, the conductance G = $10^{\mu}$ in
$G_0$, the width $\sigma$ and the FWHM (2.355 $\sigma$) in decades. The
fit is drawn in red with the conductance written next to each peak.

For every peak the text also gives the \textbf{dwell time}: how long, on
average, a trace spends in that conductance state. It comes from the
area under the fitted peak (Rashid et al.~2025, SI 2.7):

\[
\tau = \frac{A}{w} \cdot \frac{1000}{f_s}\ \mathrm{ms}, \qquad A = a\,\sigma\sqrt{2\pi}
\]

A is the area of the Gaussian (amplitude a, width $\sigma$) in the
per-trace histogram, w the bin width and f\_s the sampling rate, so A/w
is the number of samples per trace inside the peak. It needs
\emph{normalization} \texttt{per trace} or \texttt{counts}. On simulated
data with 70 \% molecular traces, 0.4 nm plateaus and 2 nm/s pulling,
the expected value is 0.7 x 0.2 s = 140 ms, and the fit gives 135-146
ms.

The program remembers the last fit; the Flicker noise tab uses it for
the peak window (9.3).

If a peak has a shoulder or a long tail, two Gaussians usually describe
it better than forcing one. And when you compare samples, compare the
per-trace histograms, not raw counts.

\hypertarget{plateau-length-and-junction-statistics}{%
\section{6. Plateau length and junction
statistics}\label{plateau-length-and-junction-statistics}}

This chapter covers two tabs: Plateau length (6.1-6.3) and Junction
statistics (6.4).

The plateau length is how far a junction can be pulled while the
molecule is still bridging the gap. Compared with the length of the
molecule it tells you whether junctions break when the molecule is fully
stretched. The tab measures it in up to three conductance windows and
also fits the tunnelling decay, which I use to calibrate the distance.

\hypertarget{settings-1}{%
\subsection{6.1 Settings}\label{settings-1}}

\begin{longtable}[]{@{}
  >{\raggedright\arraybackslash}p{(\columnwidth - 4\tabcolsep) * \real{0.2965}}
  >{\raggedright\arraybackslash}p{(\columnwidth - 4\tabcolsep) * \real{0.4727}}
  >{\raggedright\arraybackslash}p{(\columnwidth - 4\tabcolsep) * \real{0.2007}}@{}}
\toprule\noalign{}
\begin{minipage}[b]{\linewidth}\raggedright
Setting
\end{minipage} & \begin{minipage}[b]{\linewidth}\raggedright
Meaning
\end{minipage} & \begin{minipage}[b]{\linewidth}\raggedright
Default
\end{minipage} \\
\midrule\noalign{}
\endhead
\bottomrule\noalign{}
\endlastfoot
traces & All accepted traces or one cluster & all \\
window 1 from / to & Conductance window (log G), normally the peak +/-
1-2 $\sigma$ & -4.5 / -2.5 \\
window 2, window 3 & More windows, e.g.~for a second molecular state;
empty = off & off \\
length method & \texttt{span} or \texttt{count}, see below & span \\
histogram max, bins & Range and resolution of the length histogram & 2
nm, 60 \\
ignore lengths below & Leaves traces without a plateau out of the
statistics & 0 nm \\
tunnelling fit from / to & log G range for the tunnelling fit (6.3) & -5
/ -1 \\
expected $\beta$ (1/nm) & Decay constant you expect, for the calibration
in 6.3 & 20 \\
\end{longtable}

With \texttt{span}, the length runs from the first point below the upper
limit to the last point above the lower limit, so short jumps out of the
window don't make the plateau shorter. That's the usual stretching
length. \texttt{count} multiplies the number of points inside the window
by the distance step, so jumps out of the window are not counted.

\hypertarget{results}{%
\subsection{6.2 Results}\label{results}}

There's one length histogram per window, and the table gives the number
of traces, mean, median, standard deviation, skewness, kurtosis,
geometric mean, min and max for each.

Keep in mind that the real junction is longer than the measured length
by the snap-back after the gold contact breaks, roughly 0.5 nm for gold.
Add that before comparing with the molecule.

\hypertarget{tunnelling-decay-and-distance-calibration}{%
\subsection{6.3 Tunnelling decay and distance
calibration}\label{tunnelling-decay-and-distance-calibration}}

For every trace log G is fitted linearly against distance in the
tunnelling range, which gives $\beta$ in G \textasciitilde{}
exp(-$\beta$ z). The right plot is the $\beta$ histogram, the median is
written under the table.

For gold in vacuum $\beta$ is about 20 per nm (one decade of conductance
per \textasciitilde0.1 nm), in solvents it's lower; put your value in
\emph{expected beta}. The label then tells you what to multiply the
\emph{displacement ratio} (4.5) by, measured $\beta$ / expected $\beta$.
Change the ratio and detect the traces again.

\hypertarget{junction-statistics-tab}{%
\subsection{6.4 Junction statistics tab}\label{junction-statistics-tab}}

The 1D histogram tells you where the molecular peak is, but not how many
traces it came from. A peak built from 90 \% of the traces and one built
from 10 \% can look the same once you normalise per trace. The Junction
statistics tab looks at the traces one at a time and asks three
questions: did a molecule bridge the gap, for how long, and at what
conductance? From the answers you get the junction yield (the fraction
of traces with a molecular plateau, Kamenetska et al.~2009). You also
see whether the yield or the conductance changed while you were
measuring.

\textbf{How a plateau is found.} A trace ``forms a junction'' if it
stays flat inside the molecular window for at least \emph{min plateau}.
Flat means the local slope \textbar d log G / dz\textbar{} is below
\emph{max slope}. I added the slope test because without it every trace
counts. A pure tunnelling trace falls through a two-decade window in
about 0.1-0.2 nm, and with the plain \texttt{span} length from 6.1 that
already looks like a short plateau. Tunnelling drops by $\beta$ / ln 10
decades per nm, so about 9 decades/nm for gold in vacuum and 4-5 in
solvents. Molecular plateaus usually fall by well under one decade per
nm. The default of 2 decades/nm sits between the two.

The slope at each point is the change in log G between z - w/2 and z +
w/2, divided by w, where w is \emph{slope over}. Points near the ends of
a trace, where that interval doesn't fit, never count as flat.

For each trace you get:

\begin{itemize}
\tightlist
\item
  the flat length: the summed distance of all flat points inside the
  window,
\item
  the plateau conductance: the median log G of those flat points,
\item
  yes/no: does the flat length reach \emph{min plateau}?
\end{itemize}

\begin{longtable}[]{@{}
  >{\raggedright\arraybackslash}p{(\columnwidth - 4\tabcolsep) * \real{0.3164}}
  >{\raggedright\arraybackslash}p{(\columnwidth - 4\tabcolsep) * \real{0.4588}}
  >{\raggedright\arraybackslash}p{(\columnwidth - 4\tabcolsep) * \real{0.1948}}@{}}
\toprule\noalign{}
\begin{minipage}[b]{\linewidth}\raggedright
Setting
\end{minipage} & \begin{minipage}[b]{\linewidth}\raggedright
Meaning
\end{minipage} & \begin{minipage}[b]{\linewidth}\raggedright
Default
\end{minipage} \\
\midrule\noalign{}
\endhead
\bottomrule\noalign{}
\endlastfoot
traces & All accepted traces or one cluster & all \\
molecular window from / to & log G range where you expect the molecule &
-4.5 / -2.5 \\
flat if slope below & Slope limit in decades/nm & 2 \\
slope over & Distance over which the slope is measured & 0.05 nm \\
junction if flat length above & Minimum flat length for a junction & 0.1
nm \\
traces per block & Consecutive traces grouped for the plots against
trace number & 100 \\
conductance bins & Bins of the plateau conductance histogram & 60 \\
\end{longtable}

Use peak from histogram fit sets the window to the fitted peak centre
+/- 3 $\sigma$ (fit the peak in the Histograms tab first). Then click
Calculate.

\textbf{Results.} The text under the button gives the number of traces
with a plateau, the yield with its binomial standard error
sqrt(p(1-p)/N), and the median plateau conductance and length. With at
least three blocks it also gives the conductance drift in decades per
1000 traces (a straight-line fit to the block medians) and the range of
the block yields. The plots are:

\begin{longtable}[]{@{}
  >{\raggedright\arraybackslash}p{(\columnwidth - 2\tabcolsep) * \real{0.2847}}
  >{\raggedright\arraybackslash}p{(\columnwidth - 2\tabcolsep) * \real{0.6853}}@{}}
\toprule\noalign{}
\begin{minipage}[b]{\linewidth}\raggedright
Plot
\end{minipage} & \begin{minipage}[b]{\linewidth}\raggedright
What it shows
\end{minipage} \\
\midrule\noalign{}
\endhead
\bottomrule\noalign{}
\endlastfoot
Plateau conductance histogram & One entry per junction: the median log G
of its plateau. Narrower than the 1D histogram, because the tunnelling
background is gone \\
Conductance vs trace number & Every plateau conductance plus the median
of each block. A slope means the junction changed during the run
(electrode wear, solvent evaporating, the surface filling up) \\
Yield vs trace number & Yield per block with error bars and the overall
yield as a dashed line \\
Flat plateau length histogram & Flat length of every trace, with the
threshold marked. A clear gap near the threshold means the setting is
sensible \\
Length vs conductance & One point per junction. A tilt shows that
longer-stretched junctions sit at another conductance, e.g.~two binding
geometries \\
\end{longtable}

\textbf{Choosing the settings.} Look at the flat length histogram first.
Tunnelling-only traces pile up near zero and molecular traces form a
second hump. Put \emph{min plateau} in the dip between them. If there is
no dip, the yield depends on your threshold and you should say which one
you used. On simulated data (12) with 70 \% molecular traces I get 68.5
+/- 3.3 \%, and with 30 \% I get 29.0 +/- 3.2 \%.

A note on the denominator: the yield is per accepted trace. Traces
rejected in Data \& Traces (4.4) aren't counted, so loose or strict
cutting limits change it. Report the yield together with the cutting
limits.

\hypertarget{correlation-tab}{%
\section{7. Correlation tab}\label{correlation-tab}}

The correlation map shows which conductance values tend to turn up in
the same trace. That way you can see things the 1D histogram averages
away, for example a molecule going through two binding geometries one
after the other.

For each trace r the program makes its own histogram N\_r(G) on a common
set of bins, and then calculates the Pearson correlation between every
pair of bins over all traces (Makk et al.~2012):

\[
C(G_i, G_j) = \frac{\langle \delta N_r(G_i)\, \delta N_r(G_j) \rangle_r}{\sqrt{\langle \delta N_r(G_i)^2 \rangle_r \langle \delta N_r(G_j)^2 \rangle_r}}, \qquad \delta N_r = N_r - \langle N_r \rangle_r
\]

\begin{longtable}[]{@{}
  >{\raggedright\arraybackslash}p{(\columnwidth - 4\tabcolsep) * \real{0.2436}}
  >{\raggedright\arraybackslash}p{(\columnwidth - 4\tabcolsep) * \real{0.5484}}
  >{\raggedright\arraybackslash}p{(\columnwidth - 4\tabcolsep) * \real{0.1779}}@{}}
\toprule\noalign{}
\begin{minipage}[b]{\linewidth}\raggedright
Setting
\end{minipage} & \begin{minipage}[b]{\linewidth}\raggedright
Meaning
\end{minipage} & \begin{minipage}[b]{\linewidth}\raggedright
Default
\end{minipage} \\
\midrule\noalign{}
\endhead
\bottomrule\noalign{}
\endlastfoot
traces & All accepted traces or one cluster & all \\
log G min / max & Range of the map & -6 / 0.5 \\
bins / decade & Coarser than in the histograms, so each trace puts
enough counts in each bin & 20 \\
colour limit & Colour scale goes from minus to plus this & 0.3 \\
colormap & Diverging scale, red positive and blue negative & RdBu\_r \\
\end{longtable}

Click Calculate. The diagonal is always 1, and the second plot is the
mean per-trace histogram so you know where the peaks are.

How to read it: a positive region off the diagonal at (G1, G2) means
traces that stay long at G1 also stay long at G2, like two plateaus that
come together. Negative means the two exclude each other, e.g.~two
different binding geometries. Around zero there's no relation, and empty
bins have no variance and stay blank.

You need a few hundred traces. With fewer the map is mostly statistical
noise (of the order of 1/sqrt(number of traces)).

\hypertarget{clustering-tab}{%
\section{8. Clustering tab}\label{clustering-tab}}

Clustering sorts the traces into groups of similar shape without needing
a reference trace, for example molecular traces vs.~pure tunnelling, or
two binding geometries. Once the traces are clustered, the other tabs
can work on one cluster at a time.

\hypertarget{what-it-does}{%
\subsection{8.1 What it does}\label{what-it-does}}

Each accepted trace is first turned into a vector of fixed length. By
default that's a small 2D histogram (distance x log G), flattened and
normalised to sum 1, so traces of different length and sampling can be
compared (same idea as Lemmer et al.~2016 and Cabosart et al.~2019).
These vectors are centred and reduced with PCA, the clustering runs on
the first components, and finally the clusters are numbered either by
median conductance (highest first) or by size.

\hypertarget{settings-2}{%
\subsection{8.2 Settings}\label{settings-2}}

\begin{longtable}[]{@{}
  >{\raggedright\arraybackslash}p{(\columnwidth - 4\tabcolsep) * \real{0.2750}}
  >{\raggedright\arraybackslash}p{(\columnwidth - 4\tabcolsep) * \real{0.4976}}
  >{\raggedright\arraybackslash}p{(\columnwidth - 4\tabcolsep) * \real{0.1974}}@{}}
\toprule\noalign{}
\begin{minipage}[b]{\linewidth}\raggedright
Setting
\end{minipage} & \begin{minipage}[b]{\linewidth}\raggedright
Meaning
\end{minipage} & \begin{minipage}[b]{\linewidth}\raggedright
Default
\end{minipage} \\
\midrule\noalign{}
\endhead
\bottomrule\noalign{}
\endlastfoot
feature & \texttt{2D histogram} (what I'd use), \texttt{1D histogram}
(log G only) or \texttt{resampled trace} (log G on a fixed distance
grid) & 2D histogram \\
log G min / max, log G bins & Conductance grid of the feature & -6 /
0.5, 28 \\
distance min / max, distance bins & Distance grid of the feature & -0.2
/ 1.5 nm, 28 \\
resample points & Grid size for \texttt{resampled trace} & 128 \\
PCA components & How many components the clustering uses; 0 = raw
features & 10 \\
algorithm & see 8.3 & k-means \\
number of clusters & k & 3 \\
random seed & Keeps k-means, Gaussian mixture and spectral results the
same from run to run & 0 \\
order clusters by & \texttt{conductance} or \texttt{size} &
conductance \\
k range for scan & Which k the silhouette scan tries & 2, 8 \\
\end{longtable}

\hypertarget{which-algorithm}{%
\subsection{8.3 Which algorithm}\label{which-algorithm}}

\begin{longtable}[]{@{}
  >{\raggedright\arraybackslash}p{(\columnwidth - 4\tabcolsep) * \real{0.2436}}
  >{\raggedright\arraybackslash}p{(\columnwidth - 4\tabcolsep) * \real{0.4040}}
  >{\raggedright\arraybackslash}p{(\columnwidth - 4\tabcolsep) * \real{0.3223}}@{}}
\toprule\noalign{}
\begin{minipage}[b]{\linewidth}\raggedright
Algorithm
\end{minipage} & \begin{minipage}[b]{\linewidth}\raggedright
Good for
\end{minipage} & \begin{minipage}[b]{\linewidth}\raggedright
Note
\end{minipage} \\
\midrule\noalign{}
\endhead
\bottomrule\noalign{}
\endlastfoot
k-means & A first look, compact groups of similar size & Fast, assumes
round clusters \\
gaussian mixture & Groups with different spread or orientation &
Elliptical clusters, slower \\
agglomerative & Hierarchical structure & Ward linkage \\
spectral & Long or curved groups & Gets slow above \textasciitilde5000
traces \\
\end{longtable}

\hypertarget{running-it}{%
\subsection{8.4 Running it}\label{running-it}}

Run clustering uses the k you set; Run + scan k also computes the
silhouette score for each k in the scan range. The plots are the scree
plot, the PC1-PC2 scatter coloured by cluster, the 1D histogram of each
cluster, the silhouette scan (only after Run + scan k) and a 2D
histogram for every cluster. The text lists the size of each cluster and
the silhouette score (-1 to 1, higher means better separated). Clear
clusters removes the labels again.

Don't just take the k with the highest silhouette score. With noisy data
it often keeps going up with k. I take the smallest k where the cluster
2D histograms clearly look different and make physical sense; typically
that's one tunnelling cluster and one or more molecular ones.

To use a cluster, pick \texttt{cluster N} in the \emph{traces} box of
Histograms, Plateau length or Correlation. When you save results, the
histograms of every cluster are written automatically. Detecting the
traces again removes the clusters.

\hypertarget{flicker-noise-tab}{%
\section{9. Flicker noise tab}\label{flicker-noise-tab}}

The flicker noise is used to tell through-bond transport from
through-space tunnelling. The noise power NP goes with the mean
conductance as NP \textasciitilde{} $G^{n}$. If n is close to 1, the
coupling is through-bond and the current goes through the molecular
orbitals. If n is close to 2, it's through-space: electrons tunnel
directly between the electrodes, and since G \textasciitilde{}
exp(-$\beta$ d), fluctuations of the gap give dG \textasciitilde{} G and
so NP \textasciitilde{} $G^{2}$. Anything in between is a mix.

Two methods are implemented, the bell-shaped 2D histogram from Adak et
al.~(2015) and the Theil-Sen method with a stationarity test from Morris
et al.~(2025). They run on the same data and you always get both
results.

\hypertarget{what-data-you-need}{%
\subsection{9.1 What data you need}\label{what-data-you-need}}

The noise is measured on G(t) traces taken while the junction is held
still. Typically you pull until you're on the molecular plateau, stop
the piezo for 0.1-2 s, then carry on. The Data \& Traces tab finds these
holds by itself (4.3), so detect the traces first.

If your data has no holds, set \emph{segments} to
\texttt{plateau windows}; then windows inside the breaking traces are
used. These aren't as stationary as real holds, so be careful with the
result.

Morris et al.~found that n stops depending on the sampling rate above
about 5 kHz for their 90 ms traces. Shorter holds probably need faster
sampling.

\hypertarget{noise-power}{%
\subsection{9.2 Noise power}\label{noise-power}}

Each analysis window is either a whole hold (after trimming) or a
fixed-length piece of one. The first and last few milliseconds are cut
off (\emph{trim start}, \emph{trim end}) because the junction is still
relaxing right after the piezo stops. Then the PSD of G (in $G_0$) is
calculated, with Welch's method or as a periodogram. The mean of the
trace is subtracted first (\emph{detrend} \texttt{constant}); if you
don't, a trace with a step in it gives a fake 1/$f^{2}$ spectrum.
Finally the PSD is integrated over the band from \emph{integrate from}
to \emph{integrate to}:

\[
NP = \int_{f_1}^{f_2} S_G(f)\, df \qquad [G_0^2], \qquad f_1 = 100\ \mathrm{Hz},\ f_2 = 1000\ \mathrm{Hz}
\]

100-1000 Hz is the usual choice: below that there's mechanical noise,
above it amplifier noise. Whatever you choose, keep it the same for all
molecules you want to compare.

\hypertarget{which-windows-are-kept}{%
\subsection{9.3 Which windows are kept}\label{which-windows-are-kept}}

The checks below are done in this order, and every rejected window is
counted under its reason in the result text.

\begin{enumerate}
\def\labelenumi{\arabic{enumi}.}
\tightlist
\item
  Conductance range: the mean log G of the window has to be between
  \emph{accept log G from} and \emph{to}. This removes the noise floor
  and gold contacts.
\item
  Peak window: if \emph{peak centre} and \emph{peak sigma} are set, the
  mean conductance of the first and the last \emph{end fraction} (5 \%)
  of the window (or the first and last \emph{end points} points, if that
  is above 0) both have to be within centre +/- \emph{n sigma} x
  $\sigma$ (default 2). That way only junctions that stay on the
  molecular plateau are used. Peak from histogram fit copies centre and
  $\sigma$ from the last fit in the Histograms tab. If you publish, say
  which n $\sigma$ you used.
\item
  Kurtosis (optional): throws out windows with too high excess kurtosis,
  i.e.~with switching or spikes.
\item
  Drift (optional): throws out windows where the averaged log G at the
  start and at the end differ by more than the limit.
\item
  Stationarity (optional, on in the Morris preset): Augmented
  Dickey-Fuller test on log10\textbar G\textbar{} with a constant term,
  maximum lag ceil(12 (N/100)\^{}(1/4)) and the lag chosen by AIC. The
  window is kept if p \textless= \emph{ADF significance} (0.05), i.e.~if
  the unit-root hypothesis is rejected. I checked my implementation
  against statsmodels and it gives the same statistic, lag and p-value.
\end{enumerate}

With \emph{subtract instrument floor} on, the mean PSD of the windows
below \emph{floor = windows below log G} is integrated over the same
band and subtracted from every NP.

\hypertarget{how-n-is-determined}{%
\subsection{9.4 How n is determined}\label{how-n-is-determined}}

All three estimators are always calculated. \emph{estimator for n} only
decides which one is used for the bell plot and the interpretation text.

\begin{longtable}[]{@{}
  >{\raggedright\arraybackslash}p{(\columnwidth - 4\tabcolsep) * \real{0.1568}}
  >{\raggedright\arraybackslash}p{(\columnwidth - 4\tabcolsep) * \real{0.4997}}
  >{\raggedright\arraybackslash}p{(\columnwidth - 4\tabcolsep) * \real{0.3136}}@{}}
\toprule\noalign{}
\begin{minipage}[b]{\linewidth}\raggedright
Estimator
\end{minipage} & \begin{minipage}[b]{\linewidth}\raggedright
How
\end{minipage} & \begin{minipage}[b]{\linewidth}\raggedright
Comment
\end{minipage} \\
\midrule\noalign{}
\endhead
\bottomrule\noalign{}
\endlastfoot
Theil-Sen (default) & Median of the slopes between all pairs of points
in log NP vs log G, which is the same as the n where Kendall's $\tau$ of
(log G, log NP/$G^{n}$) is zero. Error from Sen's 95 \% confidence
interval & Up to 29 \% outliers don't change it. This is what Morris et
al.~recommend \\
OLS & Least-squares slope of log NP vs log G, the same as the n where
the Pearson correlation is zero & Sensitive to outliers and long tails,
tends to give n too high \\
2D Gaussian & n where the fitted 2D Gaussian (the bell) of
log(NP/$G^{n}$) vs log G has zero correlation & The graphical method of
Adak et al.~2015 \\
\end{longtable}

The OLS error is also checked with a bootstrap (\emph{bootstrap
samples}). If there are more than 6000 windows, Theil-Sen uses a random
6000 of them, otherwise it gets slow.

\hypertarget{presets}{%
\subsection{9.5 Presets}\label{presets}}

\begin{longtable}[]{@{}
  >{\raggedright\arraybackslash}p{(\columnwidth - 2\tabcolsep) * \real{0.2494}}
  >{\raggedright\arraybackslash}p{(\columnwidth - 2\tabcolsep) * \real{0.7206}}@{}}
\toprule\noalign{}
\begin{minipage}[b]{\linewidth}\raggedright
Preset
\end{minipage} & \begin{minipage}[b]{\linewidth}\raggedright
What it sets
\end{minipage} \\
\midrule\noalign{}
\endhead
\bottomrule\noalign{}
\endlastfoot
Morris et al.~2025 (ADF + Theil-Sen) & 5 ms trimmed at both ends, one
window per hold, mean-subtracted one-sided periodogram (rectangular
window), 100-1000 Hz, ADF with $\alpha$ 0.05, ends within +/- 2 $\sigma$
of the peak, no drift or kurtosis cut, Theil-Sen \\
Adak et al.~2015 (2D Gaussian bell) & 100 ms windows, Welch PSD,
100-1000 Hz, 2D Gaussian estimator \\
Rashid et al.~2025 (Pearson r scan) & First 10 ms of each hold dropped,
the rest is one window; first and last 100 points within +/- 1 $\sigma$
of the peak; DFT squared (periodogram) integrated 100-1000 Hz; no ADF,
drift or kurtosis cut; n where the Pearson r of log(NP/$G^{n}$) and log
G is smallest, scanned 0.3-2.3 in steps of 0.01 (the OLS estimator) \\
\end{longtable}

The presets don't touch your conductance range. After the Morris preset,
press Peak from histogram fit so the peak window is set.

\hypertarget{plots}{%
\subsection{9.6 Plots}\label{plots}}

\begin{longtable}[]{@{}
  >{\raggedright\arraybackslash}p{(\columnwidth - 2\tabcolsep) * \real{0.2465}}
  >{\raggedright\arraybackslash}p{(\columnwidth - 2\tabcolsep) * \real{0.7235}}@{}}
\toprule\noalign{}
\begin{minipage}[b]{\linewidth}\raggedright
Where
\end{minipage} & \begin{minipage}[b]{\linewidth}\raggedright
What
\end{minipage} \\
\midrule\noalign{}
\endhead
\bottomrule\noalign{}
\endlastfoot
Top left & The segment chosen in \emph{show segment}, accepted windows
green, rejected red \\
Top middle & Mean PSD of the accepted windows, the floor PSD, the
integration band and a 1/$f^{\alpha}$ fit \\
Top right & log NP vs log G with the OLS and Theil-Sen lines \\
Bottom left & The bell: 2D histogram of log(NP/$G^{n}$) vs log G at the
chosen n, with the fitted Gaussian as contours \\
Bottom middle & Pearson r, Kendall $\tau$ and the 2D Gaussian
correlation against n; each estimator is where its curve crosses zero \\
Bottom right & Residuals of the OLS and Theil-Sen fits; long tails mean
outliers \\
\end{longtable}

\hypertarget{what-to-report}{%
\subsection{9.7 What to report}\label{what-to-report}}

n with its error, which estimator, how many traces were accepted, the
frequency band, the trimming, the n $\sigma$ window and whether the ADF
test was on. All of this is in
\texttt{*\_\allowbreak{}Noise\_\allowbreak{}summary.txt} (13). If OLS
and Theil-Sen differ by more than their errors there are outliers in the
data; I'd trust Theil-Sen and have a look at the residual plot.

As a check I ran it on simulated data with n = 1.5 and got n\_TSE = 1.52
+/- 0.02. After adding 5 \% outlier windows, Theil-Sen gave 1.40 and OLS
1.33.

The Rashid et al.~2025 preset was checked on five simulated data sets
with 160 ms holds on the plateau (n = 1.5): it accepted 55-73 holds each
and gave n = 1.53 on average (Morris preset on the same data: 1.59). The
text also gives \emph{minimum \textbar Pearson r\textbar{} on the scan
grid}, which is how the paper reads n off the scan; it equals the OLS n
to within the grid step.

\hypertarget{events-tab-flickering-and-mechanical-switching}{%
\subsection{9.8 Events tab: flickering and mechanical
switching}\label{events-tab-flickering-and-mechanical-switching}}

A junction that switches between two states during a trace shows up as
steps in G(t): short back-and-forth steps (flickering) or one large
jump, for example when stretching converts one isomer into the other (a
mechanical event). The Events tab finds these steps the way Rashid et
al.~(2025, SI 2.6) did:

\begin{enumerate}
\def\labelenumi{\arabic{enumi}.}
\tightlist
\item
  G(t), in $G_0$, goes through a Butterworth low-pass (200 Hz, order 2)
  and then a Savitzky-Golay filter (order 20, 20 side points, so 41
  points). This removes the noise that would otherwise give false steps.
\item
  The filtered trace is differentiated with the second-order central
  difference, dG/dt = (G{[}i+1{]} - G{[}i-1{]}) / (2 dt).
\item
  Peaks of the derivative above a threshold are events: above the
  \emph{mechanical threshold} (0.002) a mechanical event, between the
  \emph{flicker threshold} (0.0007) and that a flicker.
\item
  Events where the junction is outside the molecular range (\emph{valid
  from / to log G}: metal contact, noise floor) are false positives and
  are dropped.
\end{enumerate}

Each trace is then \emph{mechanical} (at least one mechanical event),
\emph{flickering} (flickers only) or \emph{quiet}.

\begin{longtable}[]{@{}
  >{\raggedright\arraybackslash}p{(\columnwidth - 4\tabcolsep) * \real{0.3261}}
  >{\raggedright\arraybackslash}p{(\columnwidth - 4\tabcolsep) * \real{0.4314}}
  >{\raggedright\arraybackslash}p{(\columnwidth - 4\tabcolsep) * \real{0.2126}}@{}}
\toprule\noalign{}
\begin{minipage}[b]{\linewidth}\raggedright
Setting
\end{minipage} & \begin{minipage}[b]{\linewidth}\raggedright
Meaning
\end{minipage} & \begin{minipage}[b]{\linewidth}\raggedright
Default
\end{minipage} \\
\midrule\noalign{}
\endhead
\bottomrule\noalign{}
\endlastfoot
segments & \texttt{traces} (accepted breaking traces, against time) or
\texttt{holds} & traces \\
low-pass, order & Butterworth filter (applied forwards and backwards, so
no time shift) & 200 Hz, 2 \\
Savitzky-Golay side points, order & Smoothing after the low-pass & 20,
20 \\
derivative dt (s) & dt in the derivative; the thresholds are on this
scale & 0.02 \\
peak width, peak detection & See below & 50, quadratic fit \\
count & \texttt{rises}, \texttt{drops} or \texttt{both} & rises \\
flicker / mechanical threshold & Thresholds on dG/dt & 0.0007 / 0.002 \\
valid from / to log G & Molecular range & -6.5 / -1 \\
level before/after (ms) & The conductance before and after an event is
the median over this time & 20 \\
hold histogram: time bins, bins / decade & For the conductance-time
histogram of the holds & 200, 30 \\
\end{longtable}

Three details of the original LabVIEW code aren't written in the paper,
so I chose what LabVIEW does by default. Check them against your own
settings:

\begin{itemize}
\tightlist
\item
  \textbf{dt.} LabVIEW's derivative takes dt as an input. With dt = 20
  ms and 10 kHz data, the peak of dG/dt is about 2.2 times the step
  height (measured on synthetic steps), so the thresholds of 0.0007 and
  0.002 correspond to steps of about 3 x $10^{-4}$ and 9 x $10^{-4}$
  $G_0$.
\item
  \textbf{Peak width.} LabVIEW's Peak Detector fits a quadratic over
  \emph{width} points (50) and compares the fitted maximum with the
  threshold. That is \emph{peak detection} = \texttt{quadratic fit}.
  \texttt{minimum width} instead keeps only peaks at least \emph{width}
  samples wide at half height. Note that a 200 Hz low-pass makes
  derivative peaks only about 20-30 samples wide at 10 kHz, so
  \texttt{minimum width} with 50 rejects almost everything.
\item
  \textbf{Polarity.} The Peak Detector finds maxima only, so
  \emph{count} = \texttt{rises}. A flicker up and back down is then
  counted once. Use \texttt{both} to count every step.
\end{itemize}

scipy's own Savitzky-Golay coefficients fall apart at order 20 (they
come out around $10^{-17}$), so the program builds the filter from a QR
decomposition in the Legendre basis instead. It matches scipy to
$10^{-14}$ at low order and stays exact at order 20.

\textbf{Plots.} One example trace with the filtered signal and the
events marked, its derivative with both thresholds, the share of quiet,
flickering and mechanical traces, flickers per trace, and a transition
map (log G before against log G after each event). With \texttt{holds}
you also get the 2D conductance-time histogram of all holds, with the
most probable log G at every time (a Gaussian fit per time bin) and its
$\sigma$ as error bars. That is the plot the paper uses for the
mechanical modulation (Fig. S23): hold, modulate the piezo with a square
wave, and see how the most probable conductance follows.

On synthetic traces a 4 x $10^{-4}$ $G_0$ flicker at G \textasciitilde{}
$10^{-3}$ $G_0$ is found as a flicker, a jump from $10^{-6}$ to 1.6 x
$10^{-3}$ $G_0$ as a mechanical event, and a flat trace with 3 \% noise
gives nothing.

\hypertarget{i-v-egain-and-thermopower-tabs}{%
\section{10. I-V, EGaIn and Thermopower
tabs}\label{i-v-egain-and-thermopower-tabs}}

This tab is for bias sweeps. It cuts a recording into forward and
backward I-V curves, makes histograms of them and, if you want, fits a
transport model to them (10.2, 10.3).

\hypertarget{finding-the-sweeps}{%
\subsection{10.1 Finding the sweeps}\label{finding-the-sweeps}}

This first part cuts the recording into forward and backward curves and
makes histograms of them. For that the bias has to be recorded over
time, so you need a TDMS file with a bias channel or a text or Igor file
with a \emph{voltage column} (4.2). With constant bias you just get ``No
bias sweeps were found''.

The bias is smoothed and its turning points found (peaks and valleys at
least \emph{min sweep amplitude} apart). Every monotonic piece between
two turning points is one curve, forward if the bias goes up and
backward if it goes down. The current of each curve is smoothed with a
Savitzky-Golay filter and dI/dV is the ratio of the filtered derivatives
of I and V.

\begin{longtable}[]{@{}
  >{\raggedright\arraybackslash}p{(\columnwidth - 4\tabcolsep) * \real{0.3489}}
  >{\raggedright\arraybackslash}p{(\columnwidth - 4\tabcolsep) * \real{0.4518}}
  >{\raggedright\arraybackslash}p{(\columnwidth - 4\tabcolsep) * \real{0.1693}}@{}}
\toprule\noalign{}
\begin{minipage}[b]{\linewidth}\raggedright
Setting
\end{minipage} & \begin{minipage}[b]{\linewidth}\raggedright
Meaning
\end{minipage} & \begin{minipage}[b]{\linewidth}\raggedright
Default
\end{minipage} \\
\midrule\noalign{}
\endhead
\bottomrule\noalign{}
\endlastfoot
file & Which recording & first \\
min sweep amplitude (V) & Smallest bias swing that counts as a sweep &
0.05 \\
smoothing points, polynomial order & Savitzky-Golay window and order for
the current and dI/dV & 11, 2 \\
curve must reach V below / above & Only keeps curves that cover this
range; empty = off & off \\
V bins, I bins & Resolution of the 2D histograms & 100, 100 \\
log \textbar I\textbar{} histogram & Histograms log10
\textbar I\textbar{} instead of I & on \\
\end{longtable}

Top left you see up to 200 curves (forward blue, backward orange) with
the averaged forward and backward curves. The other three plots are 2D
histograms against V of log I, of log(G/$G_0$) with G = I/V, and of
log(dI/dV / $G_0$). The curves are saved as V/I column pairs in
\texttt{IV\_\allowbreak{}DATA/\allowbreak{}IV\_\allowbreak{}Curves\_\allowbreak{}F.txt}
and \texttt{IV\_\allowbreak{}Curves\_\allowbreak{}B.txt} (13).

The mean curves only use sweeps that cover the typical bias range. The
half sweeps at the start and end of a recording are left out of the
average, because they'd shrink the common range to nothing. They are
still in the histograms and in the fits of single curves.

\hypertarget{single-level-model}{%
\subsection{10.2 Single-level model}\label{single-level-model}}

Most I-V curves of single-molecule junctions can be described by one
molecular level. It sits at energy $\varepsilon_0$ from the Fermi level
of the electrodes and is broadened by its coupling $\Gamma$ to each
electrode. The transmission is a Lorentzian, and at zero temperature the
Landauer current has a closed form:

\[
I(V) = N\,G_0\,\Gamma \left[\arctan\frac{V/2 - \varepsilon(V)}{\Gamma} + \arctan\frac{V/2 + \varepsilon(V)}{\Gamma}\right], \qquad \varepsilon(V) = \varepsilon_0 + a V
\]

Energies are in eV and V in volts, so $G_0$ times an energy in eV is
directly a current in A. Both electrodes are coupled equally ($\Gamma_L$
= $\Gamma_R$ = $\Gamma$), N is the number of molecules in parallel (keep
it 1) and a describes an asymmetric junction where the level moves with
the bias. At low bias the model gives G/$G_0$ = N $\Gamma^2$ /
($\varepsilon_0^2$ + $\Gamma^2$). This is the form used, for example, by
Zotti et al.~(2010) to compare anchoring groups.

\textbf{Finite temperature.} Above 0 K the Fermi functions of the leads
smear the edges of the bias window. The general Landauer form is

\[
I = \frac{2e}{h}\int dE\, T(E)\,[f_L(E) - f_R(E)], \qquad T(E) = \frac{\Gamma^2}{(E - \varepsilon)^2 + \Gamma^2}
\]

with f\_L and f\_R the Fermi functions at $\mu$ = +V/2 and -V/2. This
T(E) is the same as 2 $\pi$ $\Gamma_L$ $\Gamma_R$ D(E) /
$\Gamma_\mathrm{tot}$, with a Lorentzian density of states D(E) of width
$\Gamma_\mathrm{tot}$ = $\Gamma_L$ + $\Gamma_R$, as written in Rashid et
al.~(2025). $\Gamma$ here is $\Gamma_{L,R}$, the number in their Table
S1. The program can do the integral in two ways:

\begin{itemize}
\tightlist
\item
  \texttt{analytic}: the integral of the Lorentzian against a Fermi
  function has a closed form with the digamma function psi (formula
  below). It is exact and fast.
\item
  \texttt{energy grid}: the integral done numerically on an energy grid,
  as in Rashid et al.~(2025, SI 2.4): -10 to 10 eV in steps of 0.02 meV.
  To keep this fast, T(E) is integrated once and the Fermi functions
  enter by one convolution (FFT) per evaluation.
\end{itemize}

The closed form used by \texttt{analytic}:

\[
\int T(E)\, f(E - \mu)\, dE = \pi\Gamma \left[\frac{1}{2} - \frac{1}{\pi}\, \mathrm{Im}\, \psi\!\left(\frac{1}{2} + \frac{\Gamma - i(\mu - \varepsilon)}{2\pi k_B T}\right)\right]
\]

The two agree to better than $10^{-9}$ (checked at 0 and 300 K,
$\varepsilon_0$ = 1.1 eV, $\Gamma$ = 11.3 meV). The single curves are
always fitted with the closed form. \texttt{energy grid} is used for the
mean and most probable curves, starting from the closed-form result, so
a grid fit takes under a second.

\textbf{Most probable curve.} Besides the mean forward and backward
curves, the program builds the most probable I-V curve. It bins all
curves by bias, fits a Gaussian to the current distribution in every bin
and takes its centre, with $\sigma$ as the error bar. This is the blue
line in Fig. S22 of Rashid et al.~(2025), and it is less sensitive to a
few odd curves than the mean. It is fitted too. \emph{weight by 1/sigma}
weights each bias bin by its spread.

\textbf{Preset: Rashid et al.~2025.} It sets 300 K, the energy grid (-10
to 10 eV, 0.02 meV), the Levenberg-Marquardt algorithm, $\Gamma_L$ =
$\Gamma_R$, no asymmetry and N = 1. The SI doesn't give the temperature,
so 300 K is my assumption; change it if your measurement was at another
one. On a model curve with the G1 values of their Table S1
($\varepsilon_0$ = 1.10 eV, $\Gamma$ = 11.3 meV, 300 K) the preset gives
back 1.100 eV and 11.30 meV.

\textbf{Ensembles.} \emph{fit N too} lets the number of molecules float,
for large-area junctions (the EGaIn tab uses it). N and $\Gamma$ are
strongly correlated then: on a model curve with N = $10^{4}$ and
$\Gamma$ = 50 meV the fit gives N = 8.6 x $10^{3}$ and $\Gamma$ = 54
meV, while $\varepsilon_0$ comes back to 0.1 \%. Report $\varepsilon_0$
and treat N and $\Gamma$ with care.

Click Fit single-level model + TVS after Analyse I-V. The fit is a
least-squares fit of I (scaled by its maximum), started from four values
of $\varepsilon_0$. The best one is kept. Every curve is fitted (up to
\emph{fit at most N single curves}), and so are the mean forward and
backward curves.

\begin{longtable}[]{@{}
  >{\raggedright\arraybackslash}p{(\columnwidth - 4\tabcolsep) * \real{0.3288}}
  >{\raggedright\arraybackslash}p{(\columnwidth - 4\tabcolsep) * \real{0.4027}}
  >{\raggedright\arraybackslash}p{(\columnwidth - 4\tabcolsep) * \real{0.2385}}@{}}
\toprule\noalign{}
\begin{minipage}[b]{\linewidth}\raggedright
Setting
\end{minipage} & \begin{minipage}[b]{\linewidth}\raggedright
Meaning
\end{minipage} & \begin{minipage}[b]{\linewidth}\raggedright
Default
\end{minipage} \\
\midrule\noalign{}
\endhead
\bottomrule\noalign{}
\endlastfoot
fit only \textbar V\textbar{} below & Restricts the fit to low bias;
empty = whole curve & empty \\
fit asymmetry a & Also fits a; off means a = 0 (symmetric junction) &
off \\
molecules in parallel N & N in the formula & 1 \\
TVS: ignore \textbar V\textbar{} below & See 10.3 & 0.05 V \\
fit at most N single curves & Fitting takes a few ms per curve & 500 \\
temperature (K) & 0 = zero-temperature closed form & 0 \\
mean-curve integration & \texttt{analytic} or \texttt{energy grid} (see
above) & analytic \\
energy grid step, from, to & Grid of the numerical integral & 0.02 meV,
-10, 10 eV \\
fit algorithm & \texttt{trust region} (bounded) or
\texttt{Levenberg-Marquardt} (unbounded) & trust region \\
fit N too & Let N float (ensembles) & off \\
most probable: bias bins, current bins & Binning for the most probable
curve & 40, 100 \\
weight by 1/$\sigma$ & Weights for the most probable curve & off \\
\end{longtable}

Some things to keep in mind:

\begin{itemize}
\tightlist
\item
  $\varepsilon_0$ comes out as a positive number. The model can't tell
  HOMO from LUMO transport, because both give the same I-V. For that you
  need thermopower or theory.
\item
  When the bias range is well below $\varepsilon_0$, the curve is almost
  linear plus a small cubic term. $\varepsilon_0$ and $\Gamma$ are then
  strongly correlated, and the errors (from the fit covariance) get
  large. Look at the $\varepsilon_0$ vs $\Gamma$ scatter: a long
  diagonal streak means the data can't separate the two.
\item
  The model assumes one level, and a level that doesn't depend on the
  bias beyond a V. At \emph{temperature} 0 it also assumes
  $\varepsilon_0$ and $\Gamma$ are much larger than kT (about 25 meV);
  otherwise set the temperature. If those assumptions don't hold, the
  model can still fit the curve nicely while $\varepsilon_0$ and
  $\Gamma$ no longer mean what they seem to.
\end{itemize}

\hypertarget{transition-voltage-spectroscopy-tvs}{%
\subsection{10.3 Transition voltage spectroscopy
(TVS)}\label{transition-voltage-spectroscopy-tvs}}

If you plot ln(|I|/$V^{2}$) against 1/V (a Fowler-Nordheim plot), I-V
curves of molecular junctions have a minimum. Beebe et al.~(2006)
introduced it as the transition voltage $V_t$. It is the same as the
maximum of $V^{2}$/|I|. $V_t$ needs no fitting and is less sensitive to
noise than the curvature of I(V), which is why it is so widely used.

For each curve and each polarity the program averages I in 60 bias bins
above \emph{TVS: ignore \textbar V\textbar{} below}, finds the maximum
of $V^{2}$/|I| and refines it with a parabola through the three highest
points. If the maximum is at the edge of the measured range, the bias
didn't go far enough, and that $V_t$ is left empty.

In the single-level model with $\Gamma$ \textless\textless{}
$\varepsilon_0$, the two transition voltages give $\varepsilon_0$ and a
in closed form (Baldea 2012):

\[
\varepsilon_0 = \frac{2\,|V_t^+ V_t^-|}{\sqrt{(V_t^+)^2 + (V_t^-)^2 + \frac{10}{3}|V_t^+ V_t^-|}}, \qquad a = \frac{(V_t^+ + V_t^-)\,\varepsilon_0}{4\,|V_t^+ V_t^-|}
\]

For a symmetric junction this reduces to $\varepsilon_0$ = (sqrt(3)/2)
$V_t$, about 0.87 $V_t$. If only one polarity has a $V_t$, that formula
is used. I checked both against the exact model: for $\varepsilon_0$ =
0.8 eV, $\Gamma$ = 10 meV and a = 0.1, the curves give $V_t^{+}$ = 1.19
V and $V_t^{-}$ = -0.75 V, and back $\varepsilon_0$ = 0.800 eV and a =
0.100.

Comparing the $\varepsilon_0$ from TVS with the $\varepsilon_0$ from the
fit is a useful check. If they disagree a lot, the curves probably
aren't single-level-like (several levels, strong asymmetry, or heating
at high bias).

\hypertarget{plots-and-what-to-report}{%
\subsection{10.4 Plots and what to
report}\label{plots-and-what-to-report}}

The fit and TVS add five plots to the I-V tab:

\begin{longtable}[]{@{}
  >{\raggedright\arraybackslash}p{(\columnwidth - 2\tabcolsep) * \real{0.3263}}
  >{\raggedright\arraybackslash}p{(\columnwidth - 2\tabcolsep) * \real{0.6437}}@{}}
\toprule\noalign{}
\begin{minipage}[b]{\linewidth}\raggedright
Plot
\end{minipage} & \begin{minipage}[b]{\linewidth}\raggedright
What it shows
\end{minipage} \\
\midrule\noalign{}
\endhead
\bottomrule\noalign{}
\endlastfoot
SLM fit & Mean forward and backward curves (dots), the most probable
curve with its $\sigma$, and the fitted model (lines) \\
Fowler-Nordheim & ln(|I|/$V^{2}$) vs 1/V of the mean curves, positive
bias solid and negative dashed, $V_t$ as dotted lines \\
Transition voltage histogram & $V_t^{+}$ and |$V_t^{-}$| of all
curves \\
SLM parameters & $\varepsilon_0$ against $\Gamma$ (log scale) for every
fitted curve \\
$\varepsilon_0$ histogram & $\varepsilon_0$ from the fits and from
$V_t$ \\
\end{longtable}

Report $\varepsilon_0$ and $\Gamma$ with their spread over the single
curves (the medians are in the text). Say whether a was fitted and give
the bias range of the fit. For TVS, give both polarities. The numbers
per curve are in
\texttt{\textless{}name\textgreater{}\_\allowbreak{}IV\_\allowbreak{}models.txt},
the most probable curve with its fit in
\texttt{\textless{}name\textgreater{}\_\allowbreak{}IV\_\allowbreak{}most\_\allowbreak{}probable.txt},
and the summary in
\texttt{\textless{}name\textgreater{}\_\allowbreak{}IV\_\allowbreak{}models\_\allowbreak{}summary.txt}
(13).

\hypertarget{egain-ensemble-tab}{%
\subsection{10.5 EGaIn (ensemble) tab}\label{egain-ensemble-tab}}

Large-area junctions (EGaIn tips, CP-AFM, crossbars) measure many
molecules at once and many junctions per sample. Their statistics differ
from break junctions in one important way: the unit is the junction, not
the sweep. A sample with 20 junctions and 20 sweeps each has 20
independent measurements, not 400. The tab follows the statistics of
Reus et al.~(2012), as implemented in GaussFit (Chiechi group) and in
the EGaIn module of XMe (Hong group).

\textbf{Loading.} With \emph{file layout} =
\texttt{one junction per file}, every file is one junction: a column
with the bias and one with the current (or J). With
\texttt{column pairs V, I, V, I, ...} one file holds many junctions:
columns 0 and 1 are the bias and current of the first junction, 2 and 3
of the second, and so on. The pairs may have different lengths (shorter
ones end in empty cells), a header line is fine, and the junctions are
named after the file with \_1, \_2, \ldots{} appended. Add files
\ldots{} or Add folder \ldots{} makes a \emph{data set}, normally one
molecule or one sample. Enter a molecular length for each data set in
the table (nm, angstrom or number of carbons, as long as all sets use
the same unit) if you want $\beta$. Add simulated makes fake data sets
with known answers.

\textbf{What it does:}

\begin{enumerate}
\def\labelenumi{\arabic{enumi}.}
\tightlist
\item
  J = I / A, with A = $\pi$ $d^{2}$ / 4 from the contact diameter, or
  the area you type in.
\item
  A junction whose \textbar J\textbar{} ever reaches the \emph{short
  threshold} is a short. It counts against the yield (working junctions
  / all junctions) and is left out.
\item
  The bias is cut into sweeps where it turns around, and every sweep is
  put on a common bias grid (\emph{bias grid step}). 0 V is left out,
  because log\textbar J\textbar{} is meaningless there.
\item
  At every bias, the log10\textbar J\textbar{} values of all sweeps are
  histogrammed (\emph{bins / decade}) and fitted with a Gaussian. The
  centre \textless log\textbar J\textbar\textgreater{} and width
  $\sigma_\mathrm{log}$ are reported, as in Reus et al.
\item
  The confidence interval uses the number of working junctions $n_j$ as
  degrees of freedom (formula below).
\item
  Rectification: sweeps are grouped into cycles that cover both
  polarities, and log R = log\textbar J(+V)\textbar{} -
  log\textbar J(-V)\textbar{} is histogrammed and fitted for every
  \textbar V\textbar.
\item
  $V_\mathrm{trans}$ (10.3) for every cycle.
\item
  Optionally the single-level model (10.2) is fitted to the
  Gaussian-mean J(V), with N free (molecules per cm$^2$ that carry the
  current).
\item
  Fit $\beta$: with two or more analysed data sets with lengths, log10
  \textless\textbar J\textbar\textgreater{} at \emph{report at bias} is
  fitted against the length d (formula below). Each point is weighted by
  the standard error of its mean, $\sigma_\mathrm{log}$ / sqrt($n_j$ -
  1). With more than two sets, the error is scaled up by the scatter of
  the points if that is larger. $\beta$ comes out per unit of the length
  you typed (Simeone et al.~2013 report 0.92 per carbon for
  alkanethiolates).
\end{enumerate}

The confidence interval (step 5) and the length dependence (step 9) are

\[
\mathrm{CI} = t_{1-\alpha/2,\; n_j - 1}\, \frac{\sigma_\mathrm{log}}{\sqrt{n_j - 1}}, \qquad \log_{10} |J| = \log_{10} J_0 - \frac{\beta\, d}{\ln 10}
\]

\begin{longtable}[]{@{}
  >{\raggedright\arraybackslash}p{(\columnwidth - 4\tabcolsep) * \real{0.3219}}
  >{\raggedright\arraybackslash}p{(\columnwidth - 4\tabcolsep) * \real{0.4176}}
  >{\raggedright\arraybackslash}p{(\columnwidth - 4\tabcolsep) * \real{0.2304}}@{}}
\toprule\noalign{}
\begin{minipage}[b]{\linewidth}\raggedright
Setting
\end{minipage} & \begin{minipage}[b]{\linewidth}\raggedright
Meaning
\end{minipage} & \begin{minipage}[b]{\linewidth}\raggedright
Default
\end{minipage} \\
\midrule\noalign{}
\endhead
\bottomrule\noalign{}
\endlastfoot
file layout & \texttt{one junction per file} or
\texttt{column pairs V, I, V, I, ...} & one junction per file \\
bias column, current column & Column number (0 = first) or header name
(one junction per file only) & 0, 1 \\
current unit & A, mA, uA, nA, or A/cm2 if the column is already J & A \\
contact diameter (um), contact area (cm2) & Geometric contact; the area
overrides the diameter & 25 um, empty \\
short if \textbar J\textbar{} reaches (A/cm2) & Short threshold & 100 \\
bias grid step (V) & Bias grid of the statistics & 0.05 \\
log\textbar J\textbar{} bins / decade & Histogram resolution & 10 \\
CI: $\alpha$ & Confidence level 1 - $\alpha$ & 0.05 \\
report at bias (V) & Bias of the log\textbar J\textbar{} and R
histograms and of $\beta$ & 0.5 \\
single-level fit, temperature & Fit of the mean J(V) & on, 300 K \\
\end{longtable}

\textbf{Plots} (for the data set chosen in \emph{show}): every sweep
with \textless log\textbar J\textbar\textgreater{} +/-
$\sigma_\mathrm{log}$, the 2D histogram of log\textbar J\textbar{}
against bias, the log\textbar J\textbar{} histogram at the report bias
with its Gaussian, the log R histogram, the $V_\mathrm{trans}$
histograms, the single-level fit and, after Fit $\beta$,
log\textbar J\textbar{} against length.

\textbf{Checks.} On simulated sets of 30 junctions with
$\sigma_\mathrm{log}$ = 0.4, a rectification log R = 0.3 and $\beta$ =
0.9, the tab finds \textless log\textbar J\textbar\textgreater{} within
its confidence interval of the truth, log R = 0.30 +/- 0.08, the shorts,
and $\beta$ = 0.904 +/- 0.016. Over 20 repeats with three sets of 20
junctions, $\beta$ came out 0.92 with a scatter of 0.087 and an average
reported error of 0.095. So the error bar is realistic: 65 \% of the
runs fall within 1 $\sigma$, against 68 \% expected.

\textbf{What to report:} \textless log\textbar J\textbar\textgreater{}
and $\sigma_\mathrm{log}$ at a stated bias, the number of junctions and
of sweeps, the yield, the contact area and how it was measured, and the
confidence interval with $\alpha$. The geometric contact area is not the
electrical one (Simeone et al.~2013 estimate $10^{-4}$ of it), so $J_0$
is only comparable between data sets measured the same way.

\hypertarget{thermopower-tab}{%
\subsection{10.6 Thermopower tab}\label{thermopower-tab}}

A temperature difference across a molecular junction drives a
thermoelectric voltage. Its sign and size give the Seebeck coefficient
S, which tells whether the HOMO or the LUMO is closer to the Fermi
level: S \textgreater{} 0 means hole transport through the HOMO, S
\textless{} 0 electron transport through the LUMO. The tab does the
standard analysis of Reddy et al.~(2007) and Widawsky et al.~(2012): you
give the thermovoltage and the temperature, and it makes the histograms
and the fit.

\textbf{Input.} Text or Igor files with the thermoelectric voltage in
one column and the temperature in another, one row per sample. A file is
usually one trace (one junction), but a file can also hold traces at
several temperatures, because the temperature is read per sample. The
temperature column is either:

\begin{itemize}
\tightlist
\item
  the temperature difference dT itself,
\item
  the temperature of the heated electrode, with dT = T - \emph{T
  reference} (the cold side, e.g.~room temperature), or
\item
  missing: then each data set gets one dT that you type in when you load
  it.
\end{itemize}

An optional conductance column (G/$G_0$ or log G/$G_0$) lets you keep
only the samples on the molecular plateau (\emph{keep log G from / to})
and plots S against log G.

\textbf{What it does:}

\begin{enumerate}
\def\labelenumi{\arabic{enumi}.}
\tightlist
\item
  V is converted to microvolt and divided by the \emph{amplifier gain}.
  \emph{invert voltage sign} flips it if your wiring measures the
  opposite polarity.
\item
  The samples are grouped by dT in steps of \emph{dT step}. With
  \emph{histogram of} = \texttt{file means}, each file contributes one
  mean voltage per dT, i.e.~one value per junction as in Reddy et
  al.~With \texttt{samples}, every sample counts.
\item
  At every dT the thermovoltage is histogrammed and fitted with a
  Gaussian: its centre is the most probable thermovoltage V(dT).
\item
  V(dT) is fitted with a straight line, weighted by the error of each
  centre. The Seebeck coefficient follows from the slope:
\end{enumerate}

\[
S = -\frac{dV}{d(\Delta T)} + S_\mathrm{electrode}
\]

\begin{enumerate}
\def\labelenumi{\arabic{enumi}.}
\setcounter{enumi}{4}
\tightlist
\item
  For every sample (or file) with dT not 0, S\_i = -V\_i/dT\_i +
  S\_electrode is histogrammed and fitted too, and plotted against log G
  if a conductance column was given.
\end{enumerate}

The error of each peak comes from resampling the files (bootstrap,
\emph{bootstrap} resamplings). The samples of one trace aren't
independent, and a Gaussian fitted to a histogram scatters more than the
plain standard error suggests. On simulated data (30 repeats, 20
junctions per dT, S = -6 uV/K) the mean S was -6.05 and 80 \% of the
results were within their error. So the error is honest, if slightly
conservative.

\emph{S of the electrodes} is added to -dV/d(dT). With 0 (the default) S
is relative to the electrode leads. Reddy et al.~add the Seebeck
coefficient of gold; enter the value for your electrode metal and
temperature if you want the absolute junction value.

\begin{longtable}[]{@{}
  >{\raggedright\arraybackslash}p{(\columnwidth - 4\tabcolsep) * \real{0.3890}}
  >{\raggedright\arraybackslash}p{(\columnwidth - 4\tabcolsep) * \real{0.4292}}
  >{\raggedright\arraybackslash}p{(\columnwidth - 4\tabcolsep) * \real{0.1518}}@{}}
\toprule\noalign{}
\begin{minipage}[b]{\linewidth}\raggedright
Setting
\end{minipage} & \begin{minipage}[b]{\linewidth}\raggedright
Meaning
\end{minipage} & \begin{minipage}[b]{\linewidth}\raggedright
Default
\end{minipage} \\
\midrule\noalign{}
\endhead
\bottomrule\noalign{}
\endlastfoot
voltage column, voltage unit & Column of the thermovoltage and its unit
(V, mV, uV, nV) & 0, uV \\
amplifier gain & V is divided by this & 1 \\
invert voltage sign & Flips the sign of V & off \\
temperature column, temperature column is & Column of the temperature
and what it holds (see above) & 1, dT \\
T reference & Cold-side temperature for the hot-side mode & 295 \\
conductance column, unit, keep log G from / to & Optional conductance
filter & empty \\
histogram of & \texttt{samples} or \texttt{file means} & samples \\
dT step & dT values are grouped to multiples of this & 1 K \\
min. samples / files per dT & Fewer than this at a dT: that dT is left
out & 20 \\
voltage bins, Seebeck bins & Histogram resolution & 80, 80 \\
S of the electrodes & Added to -dV/d(dT) & 0 uV/K \\
bootstrap & Resamplings for the peak errors & 100 \\
\end{longtable}

\textbf{Plots:} the thermovoltage histograms of all dT with their
Gaussians, the most probable V against dT with standard errors (black),
the widths (grey) and the fitted line with S, the S histogram, and the
2D histogram of S against log G.

\textbf{Checks.} On simulated traces with S = 12 and -8 uV/K, the tab
gives 11.9 +/- 0.4 and -8.1 +/- 0.4. Through real files (V in mV, the
hot-side temperature, G in $G_0$, a conductance window) S = -6 comes
back as -5.9 +/- 0.5.

\textbf{What to report:} S with its error, the dT values and the number
of junctions at each, whether samples or file means were used, the
conductance window, and whether an electrode Seebeck coefficient was
added.

\hypertarget{piezo-modulation-tab}{%
\section{11. Piezo modulation tab}\label{piezo-modulation-tab}}

If a small sine is added to the piezo, the response of the conductance
at that frequency gives the local decay constant:

\[
\beta = -\frac{d \ln G}{d z}
\]

For tunnelling $\beta$ is about 2 $\kappa$, so around 20 per nm for gold
in vacuum and less in solvents. On a molecular plateau it's small,
because stretching the junction hardly changes G.

The modulation frequency is taken as the strongest narrow line in the
piezo spectrum, unless you type it in. A line only counts if it's at
least 20 times higher than the spectrum around it, so the normal piezo
ramps aren't mistaken for a modulation. The recording is then split into
windows of a few periods; in each window ln G and z are detrended and
demodulated with a digital lock-in, and $\beta$ is minus the in-phase
ratio of the two amplitudes.

\begin{longtable}[]{@{}
  >{\raggedright\arraybackslash}p{(\columnwidth - 4\tabcolsep) * \real{0.3264}}
  >{\raggedright\arraybackslash}p{(\columnwidth - 4\tabcolsep) * \real{0.4708}}
  >{\raggedright\arraybackslash}p{(\columnwidth - 4\tabcolsep) * \real{0.1727}}@{}}
\toprule\noalign{}
\begin{minipage}[b]{\linewidth}\raggedright
Setting
\end{minipage} & \begin{minipage}[b]{\linewidth}\raggedright
Meaning
\end{minipage} & \begin{minipage}[b]{\linewidth}\raggedright
Default
\end{minipage} \\
\midrule\noalign{}
\endhead
\bottomrule\noalign{}
\endlastfoot
file & Which recording & first \\
modulation frequency (Hz) & 0 = find it automatically & 0 \\
periods per window & How long the lock-in averages & 4 \\
auto search above (Hz) & Lowest frequency the detection looks at & 10 \\
log G from / to & Only windows with mean log G in this range are shown &
-6 / 1 \\
displacement ratio & Junction nm per piezo nm (like in 4.5) & 1 \\
\end{longtable}

The top plots are the spectra of the piezo and of log G with the
modulation frequency marked, the bottom ones $\beta$ against log G for
every window and the $\beta$ histogram. Without modulation in the file
you get ``No piezo modulation was found''.

\hypertarget{simulation-tab}{%
\section{12. Simulation tab}\label{simulation-tab}}

The simulation makes a fake recording where the molecular conductance,
plateau length and noise exponent are known. I use it to get to know the
settings, to pick cutting limits and noise settings, and to check an
analysis does what it should before I trust it with real data.

Every simulated cycle goes through the same trace detection as real data
and looks like this: approach until the contact snaps in at 1-4 $G_0$;
pull through the atomic steps down to 1 $G_0$, break with snap-back, and
then (with \emph{molecule probability}) a molecular plateau, otherwise
straight tunnelling decay; hold, where G fluctuates with 1/f noise whose
power goes as $G^{n}$; retract into the noise floor; short hold at the
end.

\begin{longtable}[]{@{}
  >{\raggedright\arraybackslash}p{(\columnwidth - 4\tabcolsep) * \real{0.3589}}
  >{\raggedright\arraybackslash}p{(\columnwidth - 4\tabcolsep) * \real{0.3724}}
  >{\raggedright\arraybackslash}p{(\columnwidth - 4\tabcolsep) * \real{0.2387}}@{}}
\toprule\noalign{}
\begin{minipage}[b]{\linewidth}\raggedright
Setting
\end{minipage} & \begin{minipage}[b]{\linewidth}\raggedright
Meaning
\end{minipage} & \begin{minipage}[b]{\linewidth}\raggedright
Default
\end{minipage} \\
\midrule\noalign{}
\endhead
\bottomrule\noalign{}
\endlastfoot
number of traces & Number of cycles & 200 \\
sampling rate, speed & Hz and junction speed in nm/s & 10000, 2 \\
displacement ratio & Junction nm per piezo nm for the piezo signal &
1 \\
tunnelling decay $\beta$ & 1/nm & 10 \\
contact size range, atomic plateau length, snap-back & Contact model &
1-4 $G_0$, 0.08 nm, 0.3 nm \\
molecule probability & Fraction of traces with a molecular plateau &
0.7 \\
molecule log G, spread & Centre and spread (decades) of the molecular
conductance & -3.5, 0.25 \\
plateau length, SD, slope & Mean length and spread in nm, slope in
decades/nm & 0.4, 0.15, 0.3 \\
2nd molecule log G, fraction & A second class, to test the clustering &
off, 0 \\
hold time, pull before hold, retract after hold & s, nm, nm & 1, 1, 1 \\
noise exponent n, flicker amplitude & The real n and how big the noise
is & 1.5, 0.05 \\
noise floor log G, floor noise, white noise & Measurement floor and
white noise & -5.5, 0.3, 0.02 \\
bias, random seed & Bias in V; the same seed gives the same data & 0.1,
1 \\
\end{longtable}

Click Generate. The recording is added (or replaces what's loaded if
Replace loaded data is ticked), the traces are detected with the current
settings and the status bar shows the true mean molecular conductance
and n.~Then check: the histogram peak should be at \emph{molecule log
G}, the mean plateau length close to \emph{plateau length} (a little
off, depending on the window), the noise exponent equal to \emph{noise
exponent} within its error, and with a \emph{2nd molecule} set the
clustering should separate the two.

\hypertarget{saving-results-and-settings}{%
\section{13. Saving results and
settings}\label{saving-results-and-settings}}

File -\textgreater{} Save results (Ctrl+S) writes everything that has
been calculated into one folder, plus the figures and a complete list of
the settings that were used. You choose where the folder goes and what
it's called; by default it's the name of the first recording plus
\texttt{\_\allowbreak{}Analysis}. If the folder is already there you're
asked whether to overwrite it; No makes a new folder with
\texttt{\_\allowbreak{}2}, \texttt{\_\allowbreak{}3} and so on at the
end. Every file name starts with the folder name, written
\texttt{\textless{}name\textgreater{}} below.

The text files are tab-separated with one header line, so Origin, Excel
and Python read them directly. The figures go into a \texttt{figures}
subfolder (13.1).

\begin{longtable}[]{@{}
  >{\raggedright\arraybackslash}p{(\columnwidth - 4\tabcolsep) * \real{0.3450}}
  >{\raggedright\arraybackslash}p{(\columnwidth - 4\tabcolsep) * \real{0.3991}}
  >{\raggedright\arraybackslash}p{(\columnwidth - 4\tabcolsep) * \real{0.2259}}@{}}
\toprule\noalign{}
\begin{minipage}[b]{\linewidth}\raggedright
File
\end{minipage} & \begin{minipage}[b]{\linewidth}\raggedright
What's in it
\end{minipage} & \begin{minipage}[b]{\linewidth}\raggedright
From
\end{minipage} \\
\midrule\noalign{}
\endhead
\bottomrule\noalign{}
\endlastfoot
\texttt{\textless{}name\textgreater{}\_\allowbreak{}config.json} & All
settings of all tabs, the loaded files and the trace summary & always \\
\texttt{\textless{}name\textgreater{}\_\allowbreak{}Good\_\allowbreak{}Curves.txt}
& All accepted traces: trace, file, z\_nm, logG\_$G_0$ & Data \&
Traces \\
\texttt{\textless{}name\textgreater{}\_\allowbreak{}all\_\allowbreak{}traces.npz}
& All traces including the rejected ones (NumPy file) & Data \&
Traces \\
\texttt{\textless{}name\textgreater{}\_\allowbreak{}logHist.txt} & 1D
histogram, logG\_$G_0$ and counts & Histograms \\
\texttt{\textless{}name\textgreater{}\_\allowbreak{}3D\_\allowbreak{}Histogram.txt}
& 2D histogram as a matrix, rows = log G, columns = distance &
Histograms \\
\texttt{\textless{}name\textgreater{}\_\allowbreak{}3D\_\allowbreak{}hist\_\allowbreak{}scales.txt}
& Bin centres, distance\_nm and logG\_$G_0$ & Histograms \\
\texttt{\textless{}name\textgreater{}\_\allowbreak{}clusterN\_\allowbreak{}logHist.txt},
\texttt{...\_\allowbreak{}3D\_\allowbreak{}Histogram.txt},
\texttt{...\_\allowbreak{}3D\_\allowbreak{}hist\_\allowbreak{}scales.txt}
& The same for each cluster & Histograms, if there are clusters \\
\texttt{\textless{}name\textgreater{}\_\allowbreak{}peak\_\allowbreak{}fit.txt}
& Centre, error, G, $\sigma$, FWHM, amplitude and dwell time of each
peak & Histograms, after a fit \\
\texttt{\textless{}name\textgreater{}\_\allowbreak{}plateau length.txt}
& Plateau length of every trace, one column per window & Plateau
length \\
\texttt{\textless{}name\textgreater{}\_\allowbreak{}Plateau\_\allowbreak{}length\_\allowbreak{}parameters.txt}
& Statistics of each window and the tunnelling $\beta$ & Plateau
length \\
\texttt{\textless{}name\textgreater{}\_\allowbreak{}Correlation.txt},
\texttt{\_\allowbreak{}Correlation\_\allowbreak{}scales.txt} &
Correlation matrix; bin centres and mean histogram & Correlation \\
\texttt{\textless{}name\textgreater{}\_\allowbreak{}Clustering\_\allowbreak{}Results.txt}
& Traces per cluster, silhouette, settings & Clustering \\
\texttt{\textless{}name\textgreater{}\_\allowbreak{}cluster\_\allowbreak{}labels.txt}
& Trace, cluster, PC1-PC3 & Clustering \\
\texttt{\textless{}name\textgreater{}\_\allowbreak{}cluster\_\allowbreak{}logHist.txt}
& 1D histogram of each cluster & Clustering \\
\texttt{\textless{}name\textgreater{}\_\allowbreak{}Noise\_\allowbreak{}windows.txt}
& Every window: segment, start, samples, G, noise power, kurtosis,
drift, ADF p, accepted & Flicker noise \\
\texttt{\textless{}name\textgreater{}\_\allowbreak{}Noise\_\allowbreak{}PSD.txt}
& Frequency, mean PSD, floor PSD & Flicker noise \\
\texttt{\textless{}name\textgreater{}\_\allowbreak{}Noise\_\allowbreak{}scaling.txt}
& n, Pearson r, Kendall $\tau$, 2D Gaussian rho & Flicker noise \\
\texttt{\textless{}name\textgreater{}\_\allowbreak{}Noise\_\allowbreak{}bell\_\allowbreak{}hist.txt}
& Bell histogram as a matrix & Flicker noise \\
\texttt{\textless{}name\textgreater{}\_\allowbreak{}Noise\_\allowbreak{}summary.txt}
& n with errors for all estimators, bell parameters, counts, settings &
Flicker noise \\
\texttt{\textless{}name\textgreater{}\_\allowbreak{}IV\_\allowbreak{}DATA/\allowbreak{}IV\_\allowbreak{}Curves\_\allowbreak{}F.txt},
\texttt{IV\_\allowbreak{}Curves\_\allowbreak{}B.txt} & V/I column pairs
for every forward / backward curve & I-V \\
\texttt{\textless{}name\textgreater{}\_\allowbreak{}IV\_\allowbreak{}models.txt}
& Per curve: low-bias G, $V_t^{+}$, $V_t^{-}$, $\varepsilon_0$ and a
from TVS, $\varepsilon_0$, $\Gamma$ and a from the fit with errors,
$R^{2}$ & I-V, after the model fit \\
\texttt{\textless{}name\textgreater{}\_\allowbreak{}IV\_\allowbreak{}models\_\allowbreak{}summary.txt}
& Fits of the mean curves, medians, settings & I-V, after the model
fit \\
\texttt{\textless{}name\textgreater{}\_\allowbreak{}IV\_\allowbreak{}most\_\allowbreak{}probable.txt}
& Most probable current per bias bin, $\sigma$, points, fitted current &
I-V, after the model fit \\
\texttt{\textless{}name\textgreater{}\_\allowbreak{}events.txt} & Every
event: segment, sample, time, mechanical yes/no, dG/dt, log G before and
after & Events \\
\texttt{\textless{}name\textgreater{}\_\allowbreak{}events\_\allowbreak{}per\_\allowbreak{}trace.txt}
& Flickers and mechanical events per trace and its class & Events \\
\texttt{\textless{}name\textgreater{}\_\allowbreak{}hold\_\allowbreak{}time\_\allowbreak{}histogram.txt},
\texttt{...\_\allowbreak{}scales.txt},
\texttt{\textless{}name\textgreater{}\_\allowbreak{}hold\_\allowbreak{}most\_\allowbreak{}probable.txt}
& Conductance-time histogram of the holds and the most probable log G
per time & Events, with holds \\
\texttt{\textless{}name\textgreater{}\_\allowbreak{}events\_\allowbreak{}summary.txt}
& Counts and settings & Events \\
\texttt{\textless{}name\textgreater{}\_\allowbreak{}EGaIn\_\allowbreak{}\textless{}set\textgreater{}\_\allowbreak{}logJ.txt}
& Per bias: \textless log\textbar J\textbar\textgreater,
$\sigma_\mathrm{log}$, CI, sweeps & EGaIn \\
\texttt{\textless{}name\textgreater{}\_\allowbreak{}EGaIn\_\allowbreak{}\textless{}set\textgreater{}\_\allowbreak{}sweeps\_\allowbreak{}logJ.txt}
& log\textbar J\textbar{} of every sweep on the bias grid & EGaIn \\
\texttt{\textless{}name\textgreater{}\_\allowbreak{}EGaIn\_\allowbreak{}\textless{}set\textgreater{}\_\allowbreak{}logR.txt},
\texttt{...\_\allowbreak{}Vtrans.txt},
\texttt{...\_\allowbreak{}summary.txt} & Rectification per bias,
$V_\mathrm{trans}$ per cycle, summary & EGaIn \\
\texttt{\textless{}name\textgreater{}\_\allowbreak{}EGaIn\_\allowbreak{}beta.txt}
& $\beta$, $J_0$ and the points of the fit & EGaIn, after Fit $\beta$ \\
\texttt{\textless{}name\textgreater{}\_\allowbreak{}Thermo\_\allowbreak{}\textless{}set\textgreater{}\_\allowbreak{}peaks.txt}
& Per dT: most probable V, $\sigma$, error, count & Thermopower \\
\texttt{\textless{}name\textgreater{}\_\allowbreak{}Thermo\_\allowbreak{}\textless{}set\textgreater{}\_\allowbreak{}seebeck\_\allowbreak{}values.txt}
& S of every sample or file (and log G) & Thermopower \\
\texttt{\textless{}name\textgreater{}\_\allowbreak{}Thermo\_\allowbreak{}\textless{}set\textgreater{}\_\allowbreak{}voltage\_\allowbreak{}histograms.txt},
\texttt{...\_\allowbreak{}summary.txt} & Voltage histogram per dT; S,
slope, settings & Thermopower \\
\texttt{\textless{}name\textgreater{}\_\allowbreak{}junctions.txt} & Per
trace: junction yes/no, flat plateau length, plateau log G & Junction
statistics \\
\texttt{\textless{}name\textgreater{}\_\allowbreak{}junction\_\allowbreak{}blocks.txt}
& Per block: mean trace number, yield, its error, median plateau log G,
traces & Junction statistics \\
\texttt{\textless{}name\textgreater{}\_\allowbreak{}junction\_\allowbreak{}summary.txt}
& Yield, medians, drift, settings & Junction statistics \\
\texttt{\textless{}name\textgreater{}\_\allowbreak{}piezo\_\allowbreak{}modulation.txt}
& t, logG, amplitudes, $\beta$ and phase per window & Piezo
modulation \\
\texttt{figures/\allowbreak{}\textless{}name\textgreater{}\_\allowbreak{}\textless{}plot\textgreater{}.pdf}
and \texttt{.png} & Every plot as its own figure (13.1) & every tab \\
\texttt{\textless{}name\textgreater{}.opju} & Origin project with the
plot data and graphs, if ``plot with'' includes Origin (13.1) & every
tab \\
\end{longtable}

Plateau length, Junction statistics and Correlation results get the
cluster name added
(e.g.~\texttt{\textless{}name\textgreater{}\_\allowbreak{}cluster2\_\allowbreak{}Correlation.txt})
if they were calculated for one cluster.

\hypertarget{figures}{%
\subsection{13.1 Figures}\label{figures}}

Every plot is saved as its own figure, sized for a journal. Nothing is
combined into one image, except that 2D histograms keep their colour
bar, because you can't read them without it. Multi-panel figures can be
put together afterwards in Illustrator, Inkscape or similar.

\begin{longtable}[]{@{}
  >{\raggedright\arraybackslash}p{(\columnwidth - 4\tabcolsep) * \real{0.1517}}
  >{\raggedright\arraybackslash}p{(\columnwidth - 4\tabcolsep) * \real{0.4252}}
  >{\raggedright\arraybackslash}p{(\columnwidth - 4\tabcolsep) * \real{0.3931}}@{}}
\toprule\noalign{}
\begin{minipage}[b]{\linewidth}\raggedright
\end{minipage} & \begin{minipage}[b]{\linewidth}\raggedright
Default
\end{minipage} & \begin{minipage}[b]{\linewidth}\raggedright
Nature asks for
\end{minipage} \\
\midrule\noalign{}
\endhead
\bottomrule\noalign{}
\endlastfoot
Width & 89 mm (one column) & 89 mm single, 120 mm 1.5 columns, 183 mm
double \\
Font & Arial 7 pt, tick labels and legends 6 pt & sans-serif, 5-7 pt \\
Lines & data 0.75 pt, axes and ticks 0.5 pt & at least 0.5 pt \\
Colours & Okabe-Ito palette (colour-blind safe), viridis for maps &
colour-blind friendly \\
Titles & none, the file name says what it is & no panel titles \\
Files & PDF with editable text, plus a 600 dpi PNG & vector, or at least
300 dpi \\
\end{longtable}

All of this can be changed under File -\textgreater{} Figure export
settings (width, font, font size, line widths, tick direction, formats
pdf/png/tiff/svg/eps, resolution). The settings are remembered and also
written to \texttt{*\_\allowbreak{}config.json}.

\begin{longtable}[]{@{}
  >{\raggedright\arraybackslash}p{(\columnwidth - 4\tabcolsep) * \real{0.4936}}
  >{\raggedright\arraybackslash}p{(\columnwidth - 4\tabcolsep) * \real{0.2859}}
  >{\raggedright\arraybackslash}p{(\columnwidth - 4\tabcolsep) * \real{0.1906}}@{}}
\toprule\noalign{}
\begin{minipage}[b]{\linewidth}\raggedright
Figure files
(\texttt{figures/\allowbreak{}\textless{}name\textgreater{}\_\allowbreak{}...})
\end{minipage} & \begin{minipage}[b]{\linewidth}\raggedright
What they show
\end{minipage} & \begin{minipage}[b]{\linewidth}\raggedright
Tab
\end{minipage} \\
\midrule\noalign{}
\endhead
\bottomrule\noalign{}
\endlastfoot
\texttt{example\_\allowbreak{}traces} & Random accepted traces & Data \&
Traces \\
\texttt{logG\_\allowbreak{}histogram},
\texttt{2D\_\allowbreak{}histogram},
\texttt{clusterN\_\allowbreak{}logG\_\allowbreak{}histogram},
\texttt{clusterN\_\allowbreak{}2D\_\allowbreak{}histogram} & 1D and 2D
histograms of all traces and of each cluster & Histograms \\
\texttt{plateau\_\allowbreak{}length\_\allowbreak{}histogram},
\texttt{tunnelling\_\allowbreak{}decay\_\allowbreak{}histogram} &
Plateau lengths and tunnelling $\beta$ & Plateau length \\
\texttt{correlation\_\allowbreak{}map},
\texttt{correlation\_\allowbreak{}mean\_\allowbreak{}histogram} &
Correlation map and mean histogram & Correlation \\
\texttt{clustering\_\allowbreak{}PCA\_\allowbreak{}scree},
\texttt{clustering\_\allowbreak{}PCA\_\allowbreak{}scatter},
\texttt{clustering\_\allowbreak{}logG\_\allowbreak{}histograms},
\texttt{clustering\_\allowbreak{}silhouette\_\allowbreak{}scan},
\texttt{clustering\_\allowbreak{}clusterN\_\allowbreak{}2D\_\allowbreak{}histogram}
& PCA and cluster results & Clustering \\
\texttt{noise\_\allowbreak{}example\_\allowbreak{}segment},
\texttt{noise\_\allowbreak{}PSD},
\texttt{noise\_\allowbreak{}scaling\_\allowbreak{}fits},
\texttt{noise\_\allowbreak{}bell\_\allowbreak{}histogram},
\texttt{noise\_\allowbreak{}correlation\_\allowbreak{}vs\_\allowbreak{}n},
\texttt{noise\_\allowbreak{}fit\_\allowbreak{}residuals} & Flicker noise
& Flicker noise \\
\texttt{junction\_\allowbreak{}plateau\_\allowbreak{}conductance\_\allowbreak{}histogram},
\texttt{junction\_\allowbreak{}conductance\_\allowbreak{}vs\_\allowbreak{}trace},
\texttt{junction\_\allowbreak{}yield\_\allowbreak{}vs\_\allowbreak{}trace},
\texttt{junction\_\allowbreak{}plateau\_\allowbreak{}length\_\allowbreak{}histogram},
\texttt{junction\_\allowbreak{}length\_\allowbreak{}vs\_\allowbreak{}conductance}
& Junction yield and plateau statistics & Junction statistics \\
\texttt{IV\_\allowbreak{}curves},
\texttt{IV\_\allowbreak{}current\_\allowbreak{}histogram},
\texttt{IV\_\allowbreak{}conductance\_\allowbreak{}histogram},
\texttt{IV\_\allowbreak{}dIdV\_\allowbreak{}histogram} & I-V curves and
histograms & I-V \\
\texttt{events\_\allowbreak{}example\_\allowbreak{}trace},
\texttt{events\_\allowbreak{}example\_\allowbreak{}derivative},
\texttt{events\_\allowbreak{}trace\_\allowbreak{}classes},
\texttt{events\_\allowbreak{}flickers\_\allowbreak{}per\_\allowbreak{}trace},
\texttt{events\_\allowbreak{}transition\_\allowbreak{}map},
\texttt{events\_\allowbreak{}hold\_\allowbreak{}time\_\allowbreak{}histogram}
& Events & Events \\
\texttt{EGaIn\_\allowbreak{}\textless{}set\textgreater{}\_\allowbreak{}egain\_\allowbreak{}JV},
\texttt{...\_\allowbreak{}logJ\_\allowbreak{}2D\_\allowbreak{}histogram},
\texttt{...\_\allowbreak{}logJ\_\allowbreak{}histogram},
\texttt{...\_\allowbreak{}rectification\_\allowbreak{}histogram},
\texttt{...\_\allowbreak{}transition\_\allowbreak{}voltage\_\allowbreak{}histogram},
\texttt{...\_\allowbreak{}SLM\_\allowbreak{}fit},
\texttt{...\_\allowbreak{}beta} & Ensemble junctions, one set per data
set & EGaIn \\
\texttt{Thermo\_\allowbreak{}\textless{}set\textgreater{}\_\allowbreak{}thermo\_\allowbreak{}voltage\_\allowbreak{}histograms},
\texttt{...\_\allowbreak{}thermo\_\allowbreak{}voltage\_\allowbreak{}vs\_\allowbreak{}dT},
\texttt{...\_\allowbreak{}thermo\_\allowbreak{}seebeck\_\allowbreak{}histogram},
\texttt{...\_\allowbreak{}thermo\_\allowbreak{}seebeck\_\allowbreak{}vs\_\allowbreak{}logG}
& Thermopower, one set per data set & Thermopower \\
\texttt{IV\_\allowbreak{}SLM\_\allowbreak{}fit},
\texttt{IV\_\allowbreak{}Fowler\_\allowbreak{}Nordheim},
\texttt{IV\_\allowbreak{}transition\_\allowbreak{}voltage\_\allowbreak{}histogram},
\texttt{IV\_\allowbreak{}SLM\_\allowbreak{}parameters},
\texttt{IV\_\allowbreak{}eps0\_\allowbreak{}histogram} & Single-level
model and TVS & I-V, after the model fit \\
\texttt{modulation\_\allowbreak{}piezo\_\allowbreak{}spectrum},
\texttt{modulation\_\allowbreak{}logG\_\allowbreak{}spectrum},
\texttt{modulation\_\allowbreak{}beta\_\allowbreak{}vs\_\allowbreak{}G},
\texttt{modulation\_\allowbreak{}beta\_\allowbreak{}histogram} & Piezo
modulation & Piezo modulation \\
\end{longtable}

The plots on screen are drawn by the same code, so the exported figures
look the same, just at journal size. If you want the whole screen view
as one picture, use the save button in the plot toolbar.

\textbf{Plotting in Origin instead.} If you finish figures in Origin,
set ``plot with'' in File -\textgreater{} Figure export settings to
\texttt{Origin} (or \texttt{matplotlib + Origin} for both). Save results
then also puts every plot into Origin: the plotted data go into a
workbook (heatmaps into a matrix book) and each plot is rebuilt as an
Origin graph with the same page size, axis box, fonts, line widths,
colours, axis ranges, legend and labels as the exported figure, so they
look the same and stay editable. Histogram steps keep the original bin
centres and counts next to the step outline, and heatmaps use the same
colour map. Every tab gets its own folder in the Project Explorer, and
the project is saved as \texttt{\textless{}name\textgreater{}.opju} in
the results folder. With \texttt{Origin} alone no figure files are
written. File -\textgreater{} Send plots to Origin does the same for the
current plots without saving anything, which is handy for trying
settings.

This needs Windows, Origin 2021 or newer and the Python package
originpro (\texttt{py -m pip install originpro}). If Origin is already
open the graphs go into the open project, otherwise Origin is started;
either way it stays open. Origin is controlled from the program, so
leave it alone until the status bar says it's done (about 2 s per plot).
If Origin has a dialog open, for example a licence reminder, close it
first.

\hypertarget{settings-and-converted-files}{%
\subsection{13.2 Settings and converted
files}\label{settings-and-converted-files}}

File -\textgreater{} Save settings writes all settings to a JSON file
and File -\textgreater{} Load settings reads them back; a
\texttt{*\_\allowbreak{}config.json} from a results folder works too,
which is handy for repeating an analysis. The \texttt{.cvr} (and
\texttt{.csv}) files from the Convert tab are written separately, see
3.2.

\hypertarget{troubleshooting}{%
\section{14. Troubleshooting}\label{troubleshooting}}

\begin{longtable}[]{@{}
  >{\raggedright\arraybackslash}p{(\columnwidth - 2\tabcolsep) * \real{0.3386}}
  >{\raggedright\arraybackslash}p{(\columnwidth - 2\tabcolsep) * \real{0.6314}}@{}}
\toprule\noalign{}
\begin{minipage}[b]{\linewidth}\raggedright
Problem
\end{minipage} & \begin{minipage}[b]{\linewidth}\raggedright
What to do
\end{minipage} \\
\midrule\noalign{}
\endhead
\bottomrule\noalign{}
\endlastfoot
All traces rejected with ``never reaches high limit'' & The contact
doesn't get above \emph{high limit to check}. Lower it to just under the
contact level you see in the overview. If the amplifier saturates below
1 $G_0$, check \emph{Ia saturation} so saturated points count as contact
(3.3). \\
All traces rejected with ``never reaches low limit'' & The noise floor
is above \emph{low limit to check}. Put both low limits about half a
decade above the floor. With holds in the measurement, leave \emph{join
pull-hold-pull} on. \\
Conductance off by a constant factor & For \texttt{.cvr} files the
\emph{constant bias} in the column settings has to match the
measurement; for other files check the \emph{Ia / Ib scale} and
\emph{conductance unit}. For TDMS, look at the bias channel and series
resistor in the Convert preview. \\
Distance axis is noisy or zig-zags & More \emph{piezo smoothing}, or
\emph{distance from} \texttt{index} with the right piezo rate. \\
Distances far too long or short (MCBJ) & Calibrate the
\emph{displacement ratio} with the tunnelling fit (6.3). \\
``No hold segments were found'' & No piezo stops in the recording, or no
piezo signal at all. Use \emph{segments} = \texttt{plateau windows}, or
measure with holds. \\
``Only N accepted windows'' & Widen \emph{accept log G}, make \emph{peak
sigma} or \emph{n sigma} bigger, or switch drift, kurtosis and ADF off
one at a time to see which one throws the windows out. \\
``The integration band contains fewer than 2 frequency points'' & The
windows are too short for the band. Make \emph{window} longer or raise
\emph{integrate from}. \\
OLS and Theil-Sen give different n & There are outliers. Use Theil-Sen,
look at the residual plot, maybe switch on the ADF test. \\
``No piezo modulation was found'' & There's no modulation in the file or
it's too weak. Type in the \emph{modulation frequency} if you know
it. \\
``No bias sweeps were found'' & The bias is constant; I-V needs
sweeps. \\
No transition voltage (empty $V_t$) & The maximum of $V^{2}$/|I| is at
the end of the sweep, so the bias didn't reach $V_t$. Sweep further if
the junction survives it, or rely on the model fit. \\
SLM errors huge, $\varepsilon_0$ and $\Gamma$ on a diagonal streak & The
bias range is far below $\varepsilon_0$, so the curves can't separate
$\varepsilon_0$ and $\Gamma$. Report the low-bias conductance instead,
or measure to higher bias. \\
Events: nothing found although steps are visible & Check
\emph{derivative dt} and the thresholds: they are on the dG/dt scale
with G in $G_0$. Look at the derivative plot and put the thresholds
under the peaks you want. With \emph{peak detection} =
\texttt{minimum width}, lower \emph{peak width}. \\
Events: very many flickers on noisy data & Raise the flicker threshold
or the low-pass strength; check \emph{valid from / to log G} so the
noise floor is excluded. \\
EGaIn: every junction is a short & \emph{current unit} or \emph{contact
diameter} is wrong (J too large), or the short threshold is too low. \\
EGaIn column pairs: fewer junctions than pairs & Pairs with fewer than 4
valid points are skipped; a lone last column without a partner is
ignored. \\
Thermopower: ``fewer than \ldots{} at every dT'' & Lower \emph{min.
samples / files per dT}, or check \emph{dT step} and the temperature
column mode (dT vs hot-side temperature). \\
Thermopower: S has the wrong sign for a known molecule & Tick
\emph{invert voltage sign}; the sign of the measured voltage depends on
which electrode is grounded. \\
EGaIn: no rectification or $V_\mathrm{trans}$ & The sweeps don't cover
both polarities in one cycle, or the bias doesn't go far enough for a
$V_\mathrm{trans}$. \\
Rashid noise preset accepts almost nothing & The holds start or end off
the molecular plateau (+/- 1 $\sigma$ is strict). Check that \emph{peak
centre} is set, or relax \emph{ends within +/- n sigma}. \\
Junction yield close to 100 \% for a blank measurement & Tunnelling
counts as plateau: lower \emph{flat if slope below} or raise
\emph{junction if flat length above} (6.4). \\
Junction yield close to 0 \% although there's a clear peak & The
plateaus are tilted more than \emph{flat if slope below}, or the window
misses them; use Use peak from histogram fit. \\
Clustering changes from run to run & Keep the \emph{random seed} fixed
(it's 0 by default). \\
Window empty or tiny on a 4K screen & Set the Windows display scaling to
150-200 \%. \\
``No column \ldots{} available: \ldots{}'' & The column name or number
isn't in the file. The message lists what is there; Show columns (4.1)
does too. Numbers start at 0. \\
Igor file: ``No numeric waves found'' & The file has only text waves, or
igor2 is missing (\texttt{py -m pip install igor2}). \\
Error message about a file & The file is broken or isn't break-junction
data. Open it in a text editor or a TDMS viewer and have a look. \\
\end{longtable}

You can load several files at once; they're analysed together, the trace
numbers just continue and every trace remembers which file it came from.
The program never changes your data files, it only reads them and writes
results to the folder you pick.

To repeat an old analysis, load the \texttt{*\_\allowbreak{}config.json}
from its results folder with File -\textgreater{} Load settings, load
the same data and run the same tabs.

To check that an installation works, open a command prompt in the
program folder and run \texttt{GZero.exe --selftest report.txt} (or
\texttt{py GZero.py --selftest report.txt}). It runs the simulation,
trace detection, histogram fit, dwell time, clustering, noise analysis,
junction yield, event detection, I-V detection, the single-level fit
(also at 300 K on the energy grid with Levenberg-Marquardt) and TVS, the
EGaIn statistics with $\beta$ on simulated junctions and a column-pair
file, the thermopower fit, and the figure export without opening a
window and writes the results to \texttt{report.txt}. If everything is
fine the last line is \texttt{SELFTEST OK}.

\hypertarget{references}{%
\section{15. References}\label{references}}

I checked entries 1-24 against Crossref on 8 October 2026. The table at
the end lists what to cite for which analysis.

\hypertarget{break-junction-technique-and-background}{%
\subsection{15.1 Break-junction technique and
background}\label{break-junction-technique-and-background}}

\begin{enumerate}
\def\labelenumi{\arabic{enumi}.}
\tightlist
\item
  M. A. Reed, C. Zhou, C. J. Muller, T. P. Burgin, J. M. Tour,
  ``Conductance of a molecular junction'', \emph{Science} 278, 252-254
  (1997).
  \href{https://doi.org/10.1126/science.278.5336.252}{doi:10.1126/science.278.5336.252}
\item
  B. Xu, N. J. Tao, ``Measurement of single-molecule resistance by
  repeated formation of molecular junctions'', \emph{Science} 301,
  1221-1223 (2003).
  \href{https://doi.org/10.1126/science.1087481}{doi:10.1126/science.1087481}
\item
  L. Venkataraman, J. E. Klare, C. Nuckolls, M. S. Hybertsen, M. L.
  Steigerwald, ``Dependence of single-molecule junction conductance on
  molecular conformation'', \emph{Nature} 442, 904-907 (2006).
  \href{https://doi.org/10.1038/nature05037}{doi:10.1038/nature05037}
\item
  N. Agrait, A. Levy Yeyati, J. M. van Ruitenbeek, ``Quantum properties
  of atomic-sized conductors'', \emph{Physics Reports} 377, 81-279
  (2003).
  \href{https://doi.org/10.1016/S0370-1573(02)00633-6}{doi:10.1016/S0370-1573(02)00633-6}
\item
  T. A. Su, M. Neupane, M. L. Steigerwald, L. Venkataraman, C. Nuckolls,
  ``Chemical principles of single-molecule electronics'', \emph{Nature
  Reviews Materials} 1, 16002 (2016).
  \href{https://doi.org/10.1038/natrevmats.2016.2}{doi:10.1038/natrevmats.2016.2}
\item
  D. Xiang, X. Wang, C. Jia, T. Lee, X. Guo, ``Molecular-scale
  electronics: from concept to function'', \emph{Chemical Reviews} 116,
  4318-4440 (2016).
  \href{https://doi.org/10.1021/acs.chemrev.5b00680}{doi:10.1021/acs.chemrev.5b00680}
\end{enumerate}

\hypertarget{flicker-noise}{%
\subsection{15.2 Flicker noise}\label{flicker-noise}}

\begin{enumerate}
\def\labelenumi{\arabic{enumi}.}
\setcounter{enumi}{6}
\tightlist
\item
  O. Adak, E. Rosenthal, J. Meisner, E. F. Andrade, A. N. Pasupathy, C.
  Nuckolls, M. S. Hybertsen, L. Venkataraman, ``Flicker noise as a probe
  of electronic interaction at metal-single molecule interfaces'',
  \emph{Nano Letters} 15, 4143-4149 (2015).
  \href{https://doi.org/10.1021/acs.nanolett.5b01270}{doi:10.1021/acs.nanolett.5b01270}
  Bell-shaped NP/$G^{n}$ analysis; n = 1 through-bond, n = 2
  through-space.
\item
  A. Magyarkuti, O. Adak, A. Halbritter, L. Venkataraman, ``Electronic
  and mechanical characteristics of stacked dimer molecular junctions'',
  \emph{Nanoscale} 10, 3362-3368 (2018).
  \href{https://doi.org/10.1039/C7NR08354H}{doi:10.1039/C7NR08354H}
  Determines n from the zero of the correlation, equivalent to the OLS
  estimator in Chapter 9.4.
\item
  U. Rashid, W. Bro-Jorgensen, K. B. Harilal, P. A. Sreelakshmi, R. R.
  Mondal, V. Chittari Pisharam, K. N. Parida, K. Geetharani, J. M.
  Hamill, V. Kaliginedi, ``Chemistry of the Au-thiol interface through
  the lens of single-molecule flicker noise measurements'', \emph{J. Am.
  Chem. Soc.} 146, 9063-9073 (2024).
  \href{https://doi.org/10.1021/jacs.3c14079}{doi:10.1021/jacs.3c14079}
  A recent application of noise scaling to anchor-group chemistry.
\item
  J. M. F. Morris, J. Potter, D. Bates, C. Wu, C. M. Robertson, S. J.
  Higgins, R. J. Nichols, P. J. Low, A. Vezzoli, ``Digging through the
  (statistical) dirt: a reproducible method for single-molecule flicker
  noise analysis'', \emph{J. Phys. Chem. C} 129, 4097-4104 (2025).
  \href{https://doi.org/10.1021/acs.jpcc.4c07780}{doi:10.1021/acs.jpcc.4c07780}
  (\href{https://pmc.ncbi.nlm.nih.gov/articles/PMC11874016/}{open
  access}) ADF filter, Theil-Sen estimator, trimming. Data and code:
  \href{https://doi.org/10.17638/datacat.liverpool.ac.uk/2854}{doi:10.17638/datacat.liverpool.ac.uk/2854}.
\end{enumerate}

\hypertarget{trace-statistics-correlation-and-clustering}{%
\subsection{15.3 Trace statistics, correlation and
clustering}\label{trace-statistics-correlation-and-clustering}}

\begin{enumerate}
\def\labelenumi{\arabic{enumi}.}
\setcounter{enumi}{10}
\tightlist
\item
  P. Makk, D. Tomaszewski, J. Martinek, Z. Balogh, S. Csonka, M.
  Wawrzyniak, M. Frei, L. Venkataraman, A. Halbritter, ``Correlation
  analysis of atomic and single-molecule junction conductance'',
  \emph{ACS Nano} 6, 3411-3423 (2012).
  \href{https://doi.org/10.1021/nn300440f}{doi:10.1021/nn300440f}
\item
  M. Lemmer, M. S. Inkpen, K. Kornysheva, N. J. Long, T. Albrecht,
  ``Unsupervised vector-based classification of single-molecule charge
  transport data'', \emph{Nature Communications} 7, 12922 (2016).
  \href{https://doi.org/10.1038/ncomms12922}{doi:10.1038/ncomms12922}
\item
  D. Cabosart, M. El Abbassi, D. Stefani, R. Frisenda, M. Calame, H. S.
  J. van der Zant, M. L. Perrin, ``A reference-free clustering method
  for the analysis of molecular break-junction measurements'',
  \emph{Applied Physics Letters} 114, 143102 (2019).
  \href{https://doi.org/10.1063/1.5089198}{doi:10.1063/1.5089198} See
  also the Comment by K. P. Hoxha, M. G. Reuter, \emph{Appl. Phys.
  Lett.} 117, 136101 (2020),
  \href{https://doi.org/10.1063/5.0013531}{doi:10.1063/5.0013531}.
\item
  N. D. Bamberger, J. A. Ivie, K. N. Parida, D. V. McGrath, O. L. A.
  Monti, ``Unsupervised segmentation-based machine learning as an
  advanced analysis tool for single molecule break junction data'',
  \emph{J. Phys. Chem. C} 124, 18302-18315 (2020).
  \href{https://doi.org/10.1021/acs.jpcc.0c03612}{doi:10.1021/acs.jpcc.0c03612}
  Correction: \emph{J. Phys. Chem. C} 124, 24029-24031 (2020),
  \href{https://doi.org/10.1021/acs.jpcc.0c08721}{doi:10.1021/acs.jpcc.0c08721}.
\item
  A. Magyarkuti, N. Balogh, Z. Balogh, L. Venkataraman, A. Halbritter,
  ``Unsupervised feature recognition in single-molecule break junction
  data'', \emph{Nanoscale} 12, 8355-8363 (2020).
  \href{https://doi.org/10.1039/D0NR00467G}{doi:10.1039/D0NR00467G}
\item
  Z. Balogh, G. Mezei, N. Tenk, A. Magyarkuti, A. Halbritter,
  ``Configuration-specific insight into single-molecule conductance and
  noise data revealed by the principal component projection method'',
  \emph{J. Phys. Chem. Lett.} 14, 5109-5118 (2023).
  \href{https://doi.org/10.1021/acs.jpclett.3c00677}{doi:10.1021/acs.jpclett.3c00677}
  Combines PCA classes with noise analysis.
\end{enumerate}

\hypertarget{statistical-and-signal-processing-methods}{%
\subsection{15.4 Statistical and signal-processing
methods}\label{statistical-and-signal-processing-methods}}

\begin{enumerate}
\def\labelenumi{\arabic{enumi}.}
\setcounter{enumi}{16}
\tightlist
\item
  H. Theil, ``A rank-invariant method of linear and polynomial
  regression analysis'' (originally \emph{Proc. Koninklijke Nederlandse
  Akademie van Wetenschappen}, 1950); reprinted in \emph{Henri Theil's
  Contributions to Economics and Econometrics}, Advanced Studies in
  Theoretical and Applied Econometrics, pp.~345-381 (Springer, 1992).
  \href{https://doi.org/10.1007/978-94-011-2546-8_20}{doi:10.1007/978-94-011-2546-8\_20}
\item
  P. K. Sen, ``Estimates of the regression coefficient based on
  Kendall's $\tau$'', \emph{J. Am. Stat. Assoc.} 63, 1379-1389 (1968).
  \href{https://doi.org/10.1080/01621459.1968.10480934}{doi:10.1080/01621459.1968.10480934}
\item
  M. G. Kendall, ``A new measure of rank correlation'',
  \emph{Biometrika} 30, 81-93 (1938).
  \href{https://doi.org/10.1093/biomet/30.1-2.81}{doi:10.1093/biomet/30.1-2.81}
\item
  D. A. Dickey, W. A. Fuller, ``Distribution of the estimators for
  autoregressive time series with a unit root'', \emph{J. Am. Stat.
  Assoc.} 74, 427-431 (1979).
  \href{https://doi.org/10.1080/01621459.1979.10482531}{doi:10.1080/01621459.1979.10482531}
\item
  J. G. MacKinnon, ``Approximate asymptotic distribution functions for
  unit-root and cointegration tests'', \emph{J. Business \& Economic
  Statistics} 12, 167-176 (1994).
  \href{https://doi.org/10.1080/07350015.1994.10510005}{doi:10.1080/07350015.1994.10510005}
\item
  P. Welch, ``The use of fast Fourier transform for the estimation of
  power spectra: a method based on time averaging over short, modified
  periodograms'', \emph{IEEE Trans. Audio Electroacoust.} 15, 70-73
  (1967).
  \href{https://doi.org/10.1109/TAU.1967.1161901}{doi:10.1109/TAU.1967.1161901}
\item
  A. Savitzky, M. J. E. Golay, ``Smoothing and differentiation of data
  by simplified least squares procedures'', \emph{Analytical Chemistry}
  36, 1627-1639 (1964).
  \href{https://doi.org/10.1021/ac60214a047}{doi:10.1021/ac60214a047}
\end{enumerate}

\hypertarget{software-libraries}{%
\subsection{15.5 Software libraries}\label{software-libraries}}

\begin{enumerate}
\def\labelenumi{\arabic{enumi}.}
\setcounter{enumi}{23}
\tightlist
\item
  P. Virtanen et al., ``SciPy 1.0: fundamental algorithms for scientific
  computing in Python'', \emph{Nature Methods} 17, 261-272 (2020).
  \href{https://doi.org/10.1038/s41592-019-0686-2}{doi:10.1038/s41592-019-0686-2}
\item
  F. Pedregosa et al., ``Scikit-learn: machine learning in Python'',
  \emph{J. Machine Learning Research} 12, 2825-2830 (2011). No DOI;
  \href{https://jmlr.org/papers/v12/pedregosa11a.html}{jmlr.org}. Not in
  Crossref and not re-checked.
\end{enumerate}

\hypertarget{junction-statistics-i-v-models-events-and-ensemble-junctions}{%
\subsection{15.6 Junction statistics, I-V models, events and ensemble
junctions}\label{junction-statistics-i-v-models-events-and-ensemble-junctions}}

\begin{enumerate}
\def\labelenumi{\arabic{enumi}.}
\setcounter{enumi}{25}
\item
  M. Kamenetska, M. Koentopp, A. C. Whalley, Y. S. Park, M. L.
  Steigerwald, C. Nuckolls, M. S. Hybertsen, L. Venkataraman,
  ``Formation and evolution of single-molecule junctions'', \emph{Phys.
  Rev.~Lett.} 102, 126803 (2009).
  \href{https://doi.org/10.1103/PhysRevLett.102.126803}{doi:10.1103/PhysRevLett.102.126803}
  Junction formation probability and plateau length with molecular
  length.
\item
  L. A. Zotti, T. Kirchner, J.-C. Cuevas, F. Pauly, T. Huhn, E. Scheer,
  A. Erbe, ``Revealing the role of anchoring groups in the electrical
  conduction through single-molecule junctions'', \emph{Small} 6,
  1529-1535 (2010).
  \href{https://doi.org/10.1002/smll.200902227}{doi:10.1002/smll.200902227}
  Single-level model fits of MCBJ I-V curves.
\item
  J. M. Beebe, B. Kim, J. W. Gadzuk, C. D. Frisbie, J. G. Kushmerick,
  ``Transition from direct tunneling to field emission in
  metal-molecule-metal junctions'', \emph{Phys. Rev.~Lett.} 97, 026801
  (2006).
  \href{https://doi.org/10.1103/PhysRevLett.97.026801}{doi:10.1103/PhysRevLett.97.026801}
  Transition voltage spectroscopy.
\item
  I. Baldea, ``Ambipolar transition voltage spectroscopy: analytical
  results and experimental agreement'', \emph{Phys. Rev.~B} 85, 035442
  (2012).
  \href{https://doi.org/10.1103/PhysRevB.85.035442}{doi:10.1103/PhysRevB.85.035442}
  $\varepsilon_0$ and the bias asymmetry from $V_t^{+}$ and $V_t^{-}$.
\item
  U. Rashid, L. Medrano Sandonas, E. Chatir, Z. Ziani, P. A.
  Sreelakshmi, S. Cobo, R. Gutierrez, G. Cuniberti, V. Kaliginedi,
  ``Mapping the extended ground state reactivity landscape of a
  photoswitchable molecule at a single molecular level'', \emph{J. Am.
  Chem. Soc.} 147, 830-840 (2025).
  \href{https://doi.org/10.1021/jacs.4c13531}{doi:10.1021/jacs.4c13531}
  SI: flicker-noise protocol (2.3), single-level fit at finite
  temperature on an energy grid with Levenberg-Marquardt (2.4),
  mechanical modulation (2.5), flicker and mechanical event detection
  (2.6), dwell time (2.7).
\item
  E. M. Opodi, X. Song, X. Yu, W. Hu, ``A single level tunneling model
  for molecular junctions: evaluating the simulation methods'',
  \emph{Phys. Chem. Chem. Phys.} 24, 11958-11966 (2022).
  \href{https://doi.org/10.1039/D1CP05807J}{doi:10.1039/D1CP05807J}
  Numerical integration vs approximate formulas; see the comment by I.
  Baldea, \emph{Phys. Chem. Chem. Phys.} 26, 7230 (2024),
  \href{https://doi.org/10.1039/D2CP05110A}{doi:10.1039/D2CP05110A}.
\item
  L. J. Reus, C. A. Nijhuis, J. R. Barber, M. M. Thuo, S. Tricard, G. M.
  Whitesides, ``Statistical tools for analyzing measurements of charge
  transport'', \emph{J. Phys. Chem. C} 116, 6714-6733 (2012).
  \href{https://doi.org/10.1021/jp210445y}{doi:10.1021/jp210445y}
  Gaussian log\textbar J\textbar{} statistics and junction-based degrees
  of freedom.
\item
  F. C. Simeone, H. J. Yoon, M. M. Thuo, J. R. Barber, B. Smith, G. M.
  Whitesides, ``Defining the value of injection current and effective
  electrical contact area for EGaIn-based molecular tunneling
  junctions'', \emph{J. Am. Chem. Soc.} 135, 18131-18144 (2013).
  \href{https://doi.org/10.1021/ja408652h}{doi:10.1021/ja408652h}
  $\beta$ and $J_0$ from a length series.
\item
  Z. Pan, \ldots, W. Hong, ``XMe-Xiamen Molecular Electronics code: an
  intelligent and open-source data analysis tool for single-molecule
  conductance measurements'', \emph{Chin. J. Chem.} 42, 317 (2024).
  Code:
  \href{https://github.com/Pilab-XMU/XMe_DataAnalysis}{github.com/Pilab-XMU/XMe\_DataAnalysis}
  (Apache 2.0). Its EGaIn module (J from the contact diameter,
  log\textbar J\textbar{} statistics per bias, Gaussian fits) was used
  as a reference. DOI not checked.
\item
  GaussFit, R. C. Chiechi group:
  \href{https://github.com/rchiechi/gaussfit}{github.com/rchiechi/gaussfit}
  (GPL-3.0; now on codeberg.org/rcclab/GaussFit). EGaIn / CP-AFM
  statistics. Only its methods were used as a reference (junction-based
  confidence intervals, rectification, $V_\mathrm{trans}$ per trace). No
  code was copied, because GPL code can't go into an Apache-licensed
  program.
\item
  P. Reddy, S.-Y. Jang, R. A. Segalman, A. Majumdar, ``Thermoelectricity
  in molecular junctions'', \emph{Science} 315, 1568-1571 (2007).
  \href{https://doi.org/10.1126/science.1137149}{doi:10.1126/science.1137149}
  Thermovoltage histograms per dT and S from the slope.
\item
  J. R. Widawsky, P. Darancet, J. B. Neaton, L. Venkataraman,
  ``Simultaneous determination of conductance and thermopower of single
  molecule junctions'', \emph{Nano Letters} 12, 354-358 (2012).
  \href{https://doi.org/10.1021/nl203634m}{doi:10.1021/nl203634m}
  Conductance and thermopower in the same STM break-junction
  measurement; sign of S for HOMO- and LUMO-conducting junctions.
\end{enumerate}

I took the details of 26-37 from the publishers' and university
repositories' pages and the code repositories in October 2026. They
haven't been re-checked against Crossref like 1-24.

\hypertarget{what-to-cite-for-each-analysis}{%
\subsection{15.7 What to cite for each
analysis}\label{what-to-cite-for-each-analysis}}

\begin{longtable}[]{@{}
  >{\raggedright\arraybackslash}p{(\columnwidth - 2\tabcolsep) * \real{0.5808}}
  >{\raggedright\arraybackslash}p{(\columnwidth - 2\tabcolsep) * \real{0.3892}}@{}}
\toprule\noalign{}
\begin{minipage}[b]{\linewidth}\raggedright
Analysis in this program
\end{minipage} & \begin{minipage}[b]{\linewidth}\raggedright
Cite
\end{minipage} \\
\midrule\noalign{}
\endhead
\bottomrule\noalign{}
\endlastfoot
Conductance histograms, plateau length & 2, 3 (technique); 4 or 5
(review) \\
Junction yield, plateau conductance per trace & 26 \\
2D cross-correlation & 11 \\
Clustering & 12, 13; 15 or 14 for related methods \\
Flicker noise, bell / 2D Gaussian & 7 \\
Flicker noise, OLS / zero correlation & 8 \\
Flicker noise, ADF + Theil-Sen & 10, with 17, 18, 20 \\
Single-level model fit of I-V curves & 27 \\
Transition voltage spectroscopy & 28, 29 \\
Single-level fit at finite temperature, most probable I-V, Rashid noise
preset, events, dwell time & 30 (31 for the integration) \\
EGaIn statistics, yield, rectification & 32; 34, 35 for the software you
compare with \\
$\beta$ and $J_0$ from a length series & 33 \\
Thermopower, Seebeck coefficient & 36, 37; 34 for the XMe thermopower
module \\
PSD (Welch), smoothing & 22, 23 \\
Libraries used & 24, 25 \\
\end{longtable}

\end{document}
